\documentclass[11pt,a4paper]{article}

%─── Packages ─────────────────────────────────────────────────────────────────
\usepackage{amsmath,amssymb,amsthm}
\usepackage{geometry}
\usepackage{booktabs}
\usepackage[bookmarks=false]{hyperref}
\usepackage{array}
\usepackage{microtype}
\usepackage{hyperref}
\usepackage{xcolor}
\usepackage{enumitem}
\usepackage{graphicx}
\usepackage{float}
\geometry{top=2.5cm,bottom=2.8cm,left=2.5cm,right=2.5cm}

%─── Theorem environments ─────────────────────────────────────────────────────
\newtheorem{theorem}{Theorem}[section]
\newtheorem{lemma}[theorem]{Lemma}
\newtheorem{proposition}[theorem]{Proposition}
\newtheorem{corollary}[theorem]{Corollary}
\theoremstyle{definition}
\newtheorem{definition}[theorem]{Definition}
\newtheorem{result}[theorem]{Result}
\theoremstyle{remark}
\newtheorem{remark}[theorem]{Remark}
\newtheorem{conjecture}{Conjecture}

%─── Macros ───────────────────────────────────────────────────────────────────
\newcommand{\vp}{\varphi}
\newcommand{\ms}{\mu^{*}}
\newcommand{\lam}{\lambda}
\newcommand{\bz}{\beta_{0}}
\newcommand{\bst}{\beta^{*}}
\newcommand{\Kfb}{K_{\mathrm{fb}}}
\newcommand{\Efb}{E_{\mathrm{fb}}}
\newcommand{\Jfb}{J_{2\mathrm{iso}}^{\mathrm{fb}}}
\newcommand{\Ja}{J_{2a}}
\newcommand{\Jiso}{J_{2\mathrm{iso}}}
\newcommand{\Jfour}{J_{4}}
\newcommand{\sopt}{s_{\mathrm{opt}}}
\newcommand{\sinW}{\sin^{2}\!\theta_{W}}
\newcommand{\TCMB}{T_{\mathrm{CMB}}}
\newcommand{\rCMB}{\rho_{\mathrm{CMB}}}
\newcommand{\vEW}{v_{\mathrm{EW}}}
\newcommand{\ord}{\operatorname{ord}}
\newcommand{\leff}{\lambda_{\mathrm{eff}}}
\newcommand{\kk}{\kappa}
\newcommand{\Lcond}{\Lambda_{\mathrm{cond}}}
\newcommand{\LUV}{\Lambda_{\mathrm{UV}}}

%─── Colours ──────────────────────────────────────────────────────────────────
\definecolor{darkgreen}{RGB}{0,120,60}
\definecolor{darkred}{RGB}{160,20,20}
\definecolor{darkblue}{RGB}{0,40,140}

%══════════════════════════════════════════════════════════════════════════════
\title{\textbf{The Density-Feedback Faddeev--Niemi Hopfion:\\
Exact Algebraic Predictions for Standard-Model Parameters}
\\[0.6em]
\large Preprint --- April 2026}

\author{Frederick Manfredi}

\date{April 2026}
%══════════════════════════════════════════════════════════════════════════════

\begin{document}
\maketitle

%─── Abstract ─────────────────────────────────────────────────────────────────
\begin{abstract}
We study a density-feedback modification of the Faddeev--Niemi (FN)
Hopfion functional and show that its self-consistent saddle-point solution
carries a rich algebraic structure rooted in the McKay correspondence for
the binary icosahedral group $2I$.  Three exact algebraic identities follow
from $\vp^{2}=\vp+1$ alone, without free parameters: the feedback stiffness
$\Kfb = 2\vp\,\Ja$, the scale-selection equivalence
$\sopt = 2^{1/6} \Leftrightarrow \Jfour/\Ja = 2^{4/3}/\vp^{5}$
(Theorem~\ref{P1:thm:scale}), and the topological charge
$\mathrm{ord}(q_{2I}) = 10$.  The scale-selection value $\sopt=2^{1/6}$
is confirmed numerically to $0.010\%$ and proved in the thin-torus
approximation in the companion paper~\cite{Paper2}.
Conditionally on the $2I/E_8$ group-theoretic
input, Standard-Model observables,  $\sinW$, the electron
Yukawa $y_e$, the Koide ratio $K$, the Koide angle $\theta_K$, and
$M_W/M_Z$,  are predicted at $0.001\%$--$0.56\%$ accuracy with no
fitted parameters.  A new exact formula for the modified BPS profile,
$\Jfour/\Ja = C(2C+1)/((C^{2}+1)(3C-1))$, is derived via beta-function
integrals.  We give a precise statement of what is proven and what remains
open.
\end{abstract}

\tableofcontents
\bigskip

%══════════════════════════════════════════════════════════════════════════════
\section{Introduction}
\label{P1:sec:intro}
%══════════════════════════════════════════════════════════════════════════════

The Faddeev--Niemi (FN) Lagrangian~\cite{Faddeev1997} supports stable
topological solitons in $\mathbb{R}^3$ --- Hopfions --- characterised by a
Hopf charge $Q \in \mathbb{Z}$.  The $Q=2$ solution has toroidal geometry
and, in Battye--Sutcliffe (BS) natural units~\cite{Battye1998}, its Derrick optimal scale is determined by the ratio of two energy functionals $\Jfour/\Ja$.

In this paper we study what happens when the isotropic quadratic functional
$\Jiso$ is replaced by a \emph{density-feedback} variant
$\Jfb = \int \mathrm{kern}/(1+\bst\,\mathrm{kern})\,\mathrm{dvol}$,
motivated by a chameleon-type scalar field that mediates a
density-dependent suppression of field gradients.  The key finding is that
the modified theory possesses a self-consistent saddle point at a coupling
$\bst \approx 0.452$ at which the ratio $\Jfb/\Ja$ equals the golden ratio
$\vp$ exactly to numerical precision.  This single condition, combined with
exact phi-algebra identities, produces a cascade of Standard-Model
predictions.

\paragraph{Structure of the paper.}
Section~\ref{P1:sec:fn} defines the FN functional, the scale-invariant
saddle-finding method ($E_{\mathrm{geom}}$), and the feedback modification.
Section~\ref{P1:sec:algebra} contains the exact algebraic chain.
Section~\ref{P1:sec:group} establishes the group-theoretic origin of $\vp$
via $2I\to\hat{E}_8$, including the complete character table.
Section~\ref{P1:sec:predictions} derives all SM predictions, with the
Weinberg angle derivation in full including the $S^3$ left/right action
argument and three routes to a complete proof.
Section~\ref{P1:sec:bps} presents the modified BPS ratio theorem.
Section~\ref{P1:sec:numerics} summarises the numerical evidence.
Section~\ref{P1:sec:open} states open problems precisely.

\paragraph{What is proven vs.\ postdicted.}
We distinguish carefully throughout between
(i)~statements that are exact algebraic consequences of $\vp^{2}=\vp+1$
and the group theory of $2I$,
(ii)~numerical results confirmed on a $256\times256$ lattice to $<0.02\%$,
and
(iii)~postdictions of known SM values that follow from the framework
assuming the group-theoretic axiom.
No SM parameter has been used to fix any free parameter; the only
non-trivial input is the identification of the physical vacuum with the
$Q=2$ FN Hopfion modified by the density feedback.  Within the density-feedback
framework this identification is a consequence of the suppression formula
$S_{\mathrm{eff}}(\theta,\rho)$ and the condensate temperature, as established
in Theorem~\ref{P1:thm:vacuum}.

%══════════════════════════════════════════════════════════════════════════════
\section{The Faddeev--Niemi Functional and Its Feedback Modification}
\label{P1:sec:fn}
%══════════════════════════════════════════════════════════════════════════════

\subsection{The standard FN functional}

Work in axisymmetric coordinates $(r,z)$ with torus major radius $R_0 = 3$
(BS units).  The Hopfion profile $f : \mathbb{R}_{\geq0}\times\mathbb{R}\to[0,\pi]$
satisfies $f(0,z)=\pi$ and $f(\infty,z)=f(r,\pm\infty)=0$.  Define
\begin{equation}
  D_0 = (r-R_0)^2 + z^2, \qquad
  A   = \frac{1}{D_0} + \frac{1}{r^2}, \qquad
  \mathrm{kern} = |\nabla f|^2 + \sin^2\!f\cdot A.
  \label{P1:eq:kern}
\end{equation}
The three energy functionals are
\begin{align}
  \Ja   &= \int \sin^4\!f\cdot\mathrm{kern}\;\mathrm{dvol},
           \label{P1:eq:J2a} \\
  \Jiso &= \int \mathrm{kern}\;\mathrm{dvol},
           \label{P1:eq:Jiso} \\
  \Jfour &= \int \bigl(F_{13}^2 + F_{12}^2 + F_{23}^2\bigr)\;\mathrm{dvol},
           \label{P1:eq:J4}
\end{align}
where $\mathrm{dvol} = 2\pi r\,dr\,dz$ and
\begin{equation}
  F_{13} = -\frac{\sin^2\!f}{D_0}\bigl(f_r(r-R_0)+f_z z\bigr),
  \quad
  F_{12} = \frac{\sin^2\!f}{r}\,f_r,
  \quad
  F_{23} = \frac{\sin^2\!f}{r}\,f_z.
\end{equation}
The standard FN energy at Derrick scale~\cite{Derrick1964} $s$ is
$E(s) = s\bigl(\Ja + \ms\Jiso\bigr) + \lam \Jfour/s^3$,
optimised at $\sopt^{2} = \sqrt{\lam\Jfour/(\Ja+\ms\Jiso)}$.

Under $\mathbf{x}\mapsto s\mathbf{x}$ the integrals scale as
$\Ja\to s\Ja$, $\Jiso\to s\Jiso$, $\Jfour\to s^{-1}\Jfour$, so the
Derrick balance condition $\partial E/\partial s|_{s=1}=0$ reads
\begin{equation}
  \Ja + \ms\,\Jiso \;=\; \lam\,\Jfour.
  \label{P1:eq:derrick}
\end{equation}
This is a constraint, not a minimum: $E$ is unbounded below in the
dilation direction, and the Hopfion is therefore a \emph{saddle point}
in field-configuration space.

\subsection{The geometric-mean functional and saddle-snapshot method}
\label{P1:sec:egeom}

To locate the saddle without encountering the Derrick instability, we
use the \emph{geometric-mean functional}
\begin{equation}
  E_{\mathrm{geom}}[f;\ms]
  \;\equiv\; \bigl(\Ja + \ms\,\Jiso\bigr)\cdot\Jfour.
  \label{P1:eq:Egeom}
\end{equation}
Under $\mathbf{x}\mapsto s\mathbf{x}$,
$(\Ja+\ms\Jiso)\to s(\Ja+\ms\Jiso)$ and $\Jfour\to\Jfour/s$, so
$E_{\mathrm{geom}}$ is \emph{exactly scale-invariant}: it explores pure
shape space with the dilation mode projected out.

\begin{definition}[Saddle snapshot]
\label{P1:def:snapshot}
The \emph{saddle snapshot} is the configuration reached by gradient
descent on $E_{\mathrm{geom}}$ at the step where the Derrick scale
factor
\begin{equation}
  \sopt \;\equiv\; \sqrt{\frac{\lam\,\Jfour}{\Ja+\ms\,\Jiso}}
  \label{P1:eq:sopt_std}
\end{equation}
is closest to $1$.  At this snapshot the Derrick
balance~\eqref{P1:eq:derrick} is satisfied and the configuration is the
saddle of the full energy $E$.
\end{definition}

This method sidesteps the topological stability problem (the Hopfion
collapses to $f=0$ under unconstrained gradient descent on $E$)
by optimising the scale-invariant proxy $E_{\mathrm{geom}}$ instead.
The same snapshot method is used for the density-feedback functional
below, with $E_{\mathrm{geom}}$ replaced by $\Efb$ and $\sopt$ defined
via~\eqref{P1:eq:sopt}.

\subsection{The density-feedback modification}
\label{P1:sec:feedback}

Replace $\Jiso$ by the saturated functional
\begin{equation}
  \boxed{\Jfb \;=\; \int \frac{\mathrm{kern}}{1+\bst\,\mathrm{kern}}\;\mathrm{dvol},}
  \label{P1:eq:Jfb}
\end{equation}
motivated physically as follows.  Suppose the Hopfion vacuum is embedded
in a medium of background energy density $\rho$.  A chameleon scalar field
coupled to field gradients acquires an effective mass $m(\rho)\propto\rho$
and screens short-wavelength fluctuations.  At lowest order this produces a
suppression factor $1/(1+\bst\rho_{\mathrm{field}})$ on the gradient energy,
where $\rho_{\mathrm{field}} \propto \mathrm{kern}$.  The resulting feedback
energy is
\begin{equation}
  \Efb[f] \;=\; \Kfb\cdot\Jfour,
  \qquad
  \Kfb = \Ja + \ms\,\Jfb,
  \label{P1:eq:Efb}
\end{equation}
with Derrick scale
\begin{equation}
  \sopt^2 = \frac{\lam\,\Jfour}{\Kfb}.
  \label{P1:eq:sopt}
\end{equation}
The parameters are fixed by the group theory of $2I$
(Section~\ref{P1:sec:group}):
\begin{equation}
  \ms = 3 - \vp \;=\; 1 + \vp^{-2} \;\approx\; 1.38197,
  \qquad
  \lam = \vp^6 \;\approx\; 17.944,
  \label{P1:eq:params}
\end{equation}
where $\lam = \vp^6$ is the Bogomolny parameter uniquely determined by
the constraint $Q_{\mathrm{num}} = Q_{\mathrm{group}} = 10$
(Corollary~\ref{P1:cor:lambda}).

\begin{remark}[Physical interpretation of the feedback kernel]
\label{P1:rem:chameleon}
The functional $\Jfb = \int \mathrm{kern}/(1+\bst\,\mathrm{kern})\,\mathrm{dvol}$
admits three complementary physical interpretations.

\textbf{(i) Born--Infeld regularisation.}
$\Jfb$ is the Born--Infeld~\cite{BornInfeld1934} saturation of the isotropic FN kinetic
energy $\Jiso$: as $\mathrm{kern}\to\infty$ the integrand saturates at
$1/\bst$, preventing ultraviolet divergence of the gradient energy.
This interpretation is elevated to a derivation in
Theorem~\ref{P1:thm:vacuum} (Section~\ref{P1:sec:vacuum_thm}): $\Jfb$ is
exactly the gradient energy of the parent Lagrangian~\eqref{P1:eq:Lkin}
evaluated on the $Q=2$ Hopfion sector, making the Born--Infeld form
a consequence of the parent action rather than a postulate.

\textbf{(ii) Chameleon scalar field.}
In the chameleon mechanism~\cite{KhouryWeltman2004} a scalar field
coupled to ambient matter acquires a density-dependent effective mass
$m_{\mathrm{eff}}^2 \propto \rho_{\mathrm{matter}}$, which screens
short-wavelength fluctuations in dense regions.  Here
$\mathrm{kern}$ plays the role of local field energy density, and
$1/(1+\bst\,\mathrm{kern})$ is the corresponding screening factor.
The self-consistency condition $\Jfb/\Ja = \vp$ is the field-theoretic
fixed-point equation analogous to the chameleon equilibrium condition.

\textbf{(iii) Madelung quantum-pressure correction.}
In the weak-field expansion,
\begin{equation}
  \frac{\mathrm{kern}}{1+\bst\,\mathrm{kern}}
  \;=\; \mathrm{kern} - \bst\,\mathrm{kern}^2 + O({\bst}^2),
  \label{P1:eq:weak_field_expansion}
\end{equation}
the correction term $-\bst\,\mathrm{kern}^2$ has the same structure
as the quantum-pressure ($D^2k^4$) correction in the Madelung--Bogoliubov
dispersion relation $\omega^2 = c_L^2 k^2 + D^2 k^4$, derived from the
linearised Schr\"{o}dinger equation via the Madelung
transformation~\cite{Madelung1927,White2026}.
The feedback functional is thus a topological implementation of
Madelung hydrodynamics for an $S^2$-valued field: where
White~et~al.\ apply this dispersive correction to a scalar acoustic
medium to recover the hydrogenic spectrum, here the same correction
applied to the $\mathrm{SU}(2)_3$ WZW Hopfion selects the
self-consistent coupling $\leff^* = \lam/2^{1/3}$.
The crucial additional ingredient absent from the scalar theory is the
anisotropic $A(r,z) = D_0^{-1} + r^{-2}$ factor in $\mathrm{kern}$,
which encodes the Hopf connection geometry and selects the toroidal
Hopfion topology over spherically symmetric solutions.
\end{remark}

\begin{remark}[Topological stability of the feedback functional]
\label{P1:rem:stability}
Like the standard FN energy, $\Efb = \Kfb\cdot\Jfour$ has no
Derrick scale-barrier: both $\Kfb$ and $\Jfour$ decrease as $f\to 0$,
so the trivial vacuum $f=0$ is the global minimum of $\Efb$.  The
Hopfion is the saddle-point minimum only within the $Q=2$ topological
sector.\footnote{Throughout, $Q=2$ follows the convention of the axisymmetric
toroidal solver: the bilateral $z\to-z$ symmetry is included in the
integration measure $\mathrm{vol}=2\pi\cdot 2r\,h^2$, so the single-tube
Hopfion (which Battye--Sutcliffe~\cite{Battye1998} call $Q=1$) is labelled $Q=2$ here.
The ratio $\Kfb/\Jfour$ is independent of this bilateral factor.}  The saddle-snapshot method of Section~\ref{P1:sec:egeom}
(replacing $E$ by the scale-invariant $\Efb$ and recording the step at
which $\sopt=1$) circumvents this instability.  Numerically, the saddle
was located by a slow parameter scan from $\bst=0$,  where the
Hopfion \emph{is} a constrained local minimum and standard gradient
descent is stable,  to $\bst\approx 0.452$, warm-starting each step
from the previous converged profile.
\end{remark}

\subsection{The self-consistency condition}

\begin{definition}[Self-consistent coupling]
\label{P1:def:bstar}
$\bst$ is the unique value at which the EL minimum of $\Efb$ in the
$Q=2$ sector satisfies
\begin{equation}
  \frac{\Jfb}{\Ja} = \vp.
  \label{P1:eq:virial}
\end{equation}
Numerically, $\bst \approx 0.452$.
\end{definition}

Condition~\eqref{P1:eq:virial} plays the role of a virial balance: it is the
field-theoretic analogue of the statement that the feedback dressing of
$\Jiso$ changes by exactly the golden ratio at equilibrium.

\newpage
\begin{figure}[H]
  \centering
  \includegraphics[width=\textwidth]{hopfion_saddle_profile.png}
  \caption{%
    Saddle-point profile of the density-feedback FN Hopfion at
    $\ms=3-\vp\approx1.382$, $\bst\approx0.452$, on a
    $256\times256$ grid ($h=0.12$, $R_0=3$; BS units) plotted on a
    \emph{symmetric} $z$-axis centred at the torus ($z=0$ at the torus
    equatorial plane, $z<0$ and $z>0$ shown simultaneously).  The white
    cross marks the torus centre $(r,z)=(R_0,0)\approx(3,0)$.
    Comparing with Figure~\ref{P1:fig:profile}:
    (a) The $f=\pi/2$ Hopf tube (cyan) is closed and oval in the
    $(r,z)$-half-plane, with aspect ratio $\approx1.3$ (axially elongated).
    The oval is visibly \emph{not} round, confirming the symmetry breaking
    relative to the Battye--Sutcliffe isotropic Hopfion.
    (b,\,c) The energy density and $J_{2a}$ integrand are tightly
    concentrated inside the tube on the high-$r$ side, decaying rapidly
    outside.
    (d) Radial slice $f(\varrho)$ at $z=0$ shown \emph{without smoothing},
    making the numerical boundary-layer oscillations near
    $\varrho\approx-R_0$ explicit; these arise from the $r=0$ reflecting
    boundary on a finite grid and do not contribute to the integrated
    quantities ($J_{2a}$, $J_4$, etc.) beyond the $O(h^2)$ truncation
    error.  The BS rational ansatz (dashed) decays monotonically and does
    not cross the $f=\pi/2$ threshold at $\varrho<0$, while the physical
    profile does, confirming the topological winding.
    The header diagnostics ($\mu^*=3-\vp=1.38197$,
    $\beta_0\rho_\mathrm{CMB}=10.012$, $\mathrm{ratio}^*=0.4502$,
    $\hat{J}_{2,\mathrm{iso}}=0.5928$ vs.\ BS $0.6143$) are
    reproduced from the solver's warm-start profile without post-processing;
    the integral diagnostics ($\beta_0\rho_\mathrm{CMB}$, ratio, $s_\mathrm{opt}$)
    are from the solver's converged feedback state (Section~\ref{P1:sec:numerics}).%
  }
  \label{P1:fig:profile_sym}
\end{figure}
\newpage
\begin{figure}[H]
  \centering
  \includegraphics[width=\textwidth]{hopfion_profile_clean.png}
  \caption{%
    Post-processed complementary view of the same saddle-point solution, plotted where 
    the $z$-axis is the axial direction from the torus equatorial plane;
    the white cross marks the torus centre $(r,z)=(R_0,0)$.
    (a) Field angle $f(r,z)$: the cyan contour is the Hopf tube
    $f=\pi/2$.  The tube cross-section is visibly oval (axially elongated),
    breaking the rotational symmetry of the isotropic BS Hopfion.
    (b) Isotropic energy density $|\nabla f|^2+\sin^2\!f\cdot A$:
    the peak lies inside the Hopf tube.
    (c) Anisotropic integrand $\sin^4\!f\cdot(|\nabla f|^2+\sin^2\!f\cdot A)$:
    the $\sin^4\!f$ weighting concentrates the $J_{2a}$ contribution
    toward the tube interior, encoding the condensate anisotropy.
    (d) Radial slice at $z=0$: the physical gaussian corrected profile (red) has a higher
    peak and slower falloff than the Battye--Sutcliffe rational ansatz
    (dashed), confirming that the gradient flow has found a genuinely
    distinct shape.  The key diagnostics are
    $\beta_0\rho_{\mathrm{CMB}}=10.012$ ($0.12\%$ from $10$),
    $J_{2\mathrm{iso}}^\mathrm{fb}/J_{2a}=1.6182\approx\vp$
    ($0.012\%$), and $s_\mathrm{opt}=2^{1/6}$ to $0.010\%$,
    as reported by the solver on the converged feedback profile
    (Section~\ref{P1:sec:numerics}); the displayed profile is the 
    warm-start used to seed the feedback gradient descent.%
  }
  \label{P1:fig:profile}
\end{figure}

\newpage
%══════════════════════════════════════════════════════════════════════════════
\subsection{A radial Pohozaev identity: the $P_{\Kfb}$ algebraic result}
\label{P1:sec:pohozaev2}
%══════════════════════════════════════════════════════════════════════════════

The Derrick balance (first Pohozaev identity, Section~\ref{P1:sec:fn}
of this paper, the \emph{isotropic} dilation $\xi=rf_r+zf_z$)
contracts the EL equation with the isotropic dilation
$\xi = r\partial_r f + z\partial_z f$.
A \emph{radial} Pohozaev identity contracts instead with
$\xi = r\partial_r f$ (holding $z$ and $R_0$ fixed), generating a
topological boundary term from the core shift.
The full development of this identity, including the linear equation
for $\Ja$ and its structural implications, is recorded in Appendix A of Paper~III~\cite{Paper3}, Variation (d) therein, where it joins the three earlier
Pohozaev variations that also supported but did not close the main proof.

One result from that analysis is a clean algebraic identity that
belongs here because it is a direct consequence of the Derrick balance
already proved in this paper.

\begin{proposition}[$P_{\Kfb} = 2\vp\,P_{\Ja}$ is an algebraic identity]
\label{P1:prop:PKfb}
Let $P_I = d I[f(r(1+\varepsilon),z)]/d\varepsilon|_0$ denote the
variation of functional $I$ under the radial dilation
$r\mapsto r(1+\varepsilon)$, $z$ and $R_0$ fixed.
At the self-consistent coupling $\bst$ satisfying
$\Kfb = \Ja + \ms\Jfb = 2\vp\Ja$ and $\ms = 3-\vp$:
\begin{equation}
  P_{\Kfb} \;=\; 2\vp\,P_{\Ja}.
  \label{P1:eq:PKfb_exact}
\end{equation}
\end{proposition}

\begin{proof}
$P_{\Kfb} = P_{\Ja} + \ms P_{\Jfb}$.
For~\eqref{P1:eq:PKfb_exact} we need $P_{\Jfb}/P_{\Ja} = (2\vp-1)/\ms
= (2\vp-1)/(3-\vp)$.
This ratio equals $\vp$ identically:
$(2\vp-1)/(3-\vp)=\vp \Leftrightarrow 2\vp-1=3\vp-\vp^2=3\vp-(\vp+1)=2\vp-1$,
where we used $\vp^2=\vp+1$.
Substituting back: $P_{\Kfb}=P_{\Ja}(1+(3-\vp)\vp)=P_{\Ja}(1+3\vp-\vp-1)=2\vp P_{\Ja}$.
\qed
\end{proof}

\newpage
%══════════════════════════════════════════════════════════════════════════════
\section{The Exact Algebraic Chain}
\label{P1:sec:algebra}
%══════════════════════════════════════════════════════════════════════════════

Everything in this section is exact, following purely from
$\vp^{2}=\vp+1$ and Definition~\ref{P1:def:bstar}.  No numerical
approximation is used.

\begin{lemma}[Golden identity]
\label{P1:lem:golden}
$1 + \ms\vp = 2\vp$ exactly.
\end{lemma}

\begin{proof}
$1 + (3-\vp)\vp = 1 + 3\vp - \vp^{2} = 1 + 3\vp - (\vp+1) = 2\vp$.
\end{proof}

\begin{lemma}[Virial sum]
\label{P1:lem:virial}
$\vp^{-1} + \ms = 2$ exactly.
\end{lemma}

\begin{proof}
$\vp^{-1} = \vp - 1$, so $(\vp-1)+(3-\vp) = 2$.
\end{proof}

\begin{proposition}[Feedback stiffness]
\label{P1:prop:Kfb}
At the self-consistent $\bst$,
$\Kfb = 2\vp\,\Ja$.
\end{proposition}

\begin{proof}
$\Kfb = \Ja + \ms\Jfb = \Ja(1 + \ms\cdot\Jfb/\Ja) = \Ja(1+\ms\vp) = 2\vp\,\Ja$
by~\eqref{P1:eq:virial} and Lemma~\ref{P1:lem:golden}.
\end{proof}

\begin{theorem}[Scale selection]
\label{P1:thm:scale}
At the self-consistent $\bst$,
\begin{equation}
  \sopt = 2^{1/6}
  \quad\Longleftrightarrow\quad
  \frac{\Jfour}{\Ja} = \frac{2^{4/3}}{\vp^{5}}.
  \label{P1:eq:equiv}
\end{equation}
These two conditions are algebraically identical; neither is more primitive
than the other.
\end{theorem}

\begin{proof}
From~\eqref{P1:eq:sopt} and Proposition~\ref{P1:prop:Kfb}:
\[
  \sopt^{2} = \frac{\vp^{6}\Jfour}{2\vp\,\Ja} = \frac{\vp^{5}}{2}\cdot\frac{\Jfour}{\Ja}.
\]
Then $\sopt=2^{1/6}\Leftrightarrow\sopt^{6}=2\Leftrightarrow
(\vp^{5}/2)^{3}(\Jfour/\Ja)^{3}=2\Leftrightarrow\Jfour/\Ja = (16/\vp^{15})^{1/3}
= 2^{4/3}/\vp^{5}$.
\end{proof}

\begin{remark}[Connection to $E_8$ and WZW]
\label{P1:rem:e8}
The value $\sopt = 2^{1/(2k)}$ with $k=3$ follows from the Coxeter number
of $E_8$: $h(E_8) = 30 = 2k(k+2)$ is satisfied uniquely by $k=3$ among
ADE algebras with $h\leq 30$.  The WZW level $k=3$ is forced by
$\mathbb{Z}_{10}\subset 2I$ (Section~\ref{P1:sec:group}).  The factor of~$2$
in $\sopt^{6}=2$ is $|2I|/|I|=2$, the order of the central extension.
\end{remark}

%══════════════════════════════════════════════════════════════════════════════
\section{Group-Theoretic Origin of $\varphi$}
\label{P1:sec:group}
%══════════════════════════════════════════════════════════════════════════════

\subsection{The binary icosahedral group and $\hat{E}_8$}

The \emph{binary icosahedral group} $2I \subset \mathrm{SU}(2)$ has order
$|2I|=120=2|I|$ and nine irreducible representations of dimensions
$1,2,2,3,3,4,4,5,6$.  The golden ratio appears as the character value
$\chi_{\Gamma_{2}}(C_{5}) = \vp$ in the two-dimensional irrep $\Gamma_2$.

The nine conjugacy classes, with their sizes, orders in $2I$, and
$\mathrm{SU}(2)$ rotation angles, are:
\begin{equation}
\renewcommand{\arraystretch}{1.25}
\begin{array}{lrrl}
\toprule
\text{Class} & |C_k| & \text{ord} & \text{SU}(2)\text{ angle}\\
\midrule
1       & 1  & 1  & 0 \\
\bar{1} & 1  & 2  & 2\pi \\
C_5     & 12 & 10 & 2\pi/5 \\
C_5^2   & 12 & 5  & 4\pi/5 \\
C_5^3   & 12 & 10 & 6\pi/5 \\
C_5^4   & 12 & 5  & 8\pi/5 \\
C_3     & 20 & 6  & 2\pi/3 \\
C_3^2   & 20 & 3  & 4\pi/3 \\
C_2     & 30 & 4  & \pi    \\
\bottomrule
\end{array}
\label{P1:eq:classes}
\end{equation}
(Class sizes sum to $120$; orders follow from $n\theta\equiv 0\pmod{4\pi}$
in $\mathrm{SU}(2)$.)  Setting $\bar\vp = 1/\vp = \vp-1$, the complete
character table is:
\begin{equation}
\renewcommand{\arraystretch}{1.3}
\begin{array}{l|rrrrrrrrr}
\toprule
\Gamma & 1 & \bar{1} & C_5 & C_5^2 & C_5^3 & C_5^4 & C_3 & C_3^2 & C_2\\
\midrule
\Gamma_1\;[1]  & 1 & 1  & 1         & 1          & 1           & 1         & 1  &  1 &  1\\
\Gamma_2\;[2]  & 2 &-2  & \vp       & \bar\vp    & -\bar\vp    & -\vp      & 1  & -1 &  0\\
\Gamma_3\;[2'] & 2 &-2  & -\bar\vp  & -\vp       & \vp         & \bar\vp   & 1  & -1 &  0\\
\Gamma_4\;[3]  & 3 & 3  & \vp       & -\bar\vp   & -\bar\vp    & \vp       & 0  &  0 & -1\\
\Gamma_5\;[3'] & 3 & 3  & -\bar\vp  & \vp        & \vp         & -\bar\vp  & 0  &  0 & -1\\
\Gamma_6\;[4]  & 4 & 4  & -1        & -1         & -1          & -1        & 1  &  1 &  0\\
\Gamma_7\;[4'] & 4 &-4  & 1         & -1         & 1           & -1        & -1 &  1 &  0\\
\Gamma_8\;[5]  & 5 & 5  & 0         & 0          & 0           & 0         & -1 & -1 &  1\\
\Gamma_9\;[6]  & 6 &-6  & -1        & 1          & -1          & 1         & 0  &  0 &  0\\
\bottomrule
\end{array}
\label{P1:eq:chartable}
\end{equation}
The values $\chi_{\Gamma_2}(C_5)=\vp$ and $\chi_{\Gamma_2}(C_5^2)=\bar\vp$ follow from the $\mathrm{SU}(2)$ spin-$\tfrac{1}{2}$ formula $\chi_{1/2}(\theta)=2\cos(\theta/2)$ at $\theta=2\pi/5$ and $4\pi/5$ respectively.\\
All 81 column-orthogonality relations $\sum_k|C_k|\chi_i(C_k)\overline{\chi_j(C_k)} = 120\,\delta_{ij}$ are satisfied identically.

\begin{theorem}[McKay correspondence~\cite{McKay1980}]
\label{P1:thm:mckay}
The McKay graph of $2I$ constructed from $\Gamma_2$ is the affine Dynkin
diagram $\hat{E}_8$, establishing the chain
\begin{equation}
  2I \;\longleftrightarrow\; \hat{E}_8 \;\supset\;
  E_8 \;\supset\; \mathrm{SU}(5) \;\supset\;
  \mathrm{SU}(3)_C\times\mathrm{SU}(2)_L\times U(1)_Y.
  \label{P1:eq:chain}
\end{equation}
\end{theorem}

\subsection{Forcing the WZW level}

\begin{theorem}[WZW level $k=3$]
\label{P1:thm:k3}
The subgroup $\mathbb{Z}_{10}\subset 2I$ forces the $\mathrm{SU}(2)$ WZW
model to level $k=3$ uniquely via
\begin{equation}
  \ord(q) = 2(k+2)\big|_{k=3} = 10 = |\mathbb{Z}_{10}|,
  \qquad
  q = e^{i\pi/(k+2)} = e^{i\pi/5}.
  \label{P1:eq:q}
\end{equation}
\end{theorem}

\subsection{$\varphi$ from the quantum group}

\begin{theorem}[Quantum dimension equals $\vp$]
\label{P1:thm:dimq}
The quantum dimension of the spin-$\tfrac{1}{2}$ primary at WZW level $k=3$
equals $\vp$:
\begin{equation}
  \dim_q\!\left(\tfrac{1}{2};\,k=3\right)
  = \frac{S_{0,\,1/2}}{S_{0,0}}
  = \frac{\sin(2\pi/5)}{\sin(\pi/5)}
  = 2\cos(\pi/5)
  = \vp,
  \label{P1:eq:dimq}
\end{equation}
where $S$ is the modular $S$-matrix of $\mathrm{SU}(2)_3$~\cite{Verlinde1988}.  This
coincides with the character value $\chi_{\Gamma_2}(C_5)=\vp$ of the
$2I$ representation and is the unique fixed point of the recursion
$x^2 = x+1$.
\end{theorem}

The self-consistency condition $\Jfb/\Ja=\vp$ (Definition~\ref{P1:def:bstar})
therefore has a direct group-theoretic interpretation: the feedback ratio
at equilibrium equals the quantum dimension of the fundamental
representation of the $\mathrm{SU}(2)_3$ WZW model, which is itself
determined by the icosahedral symmetry of the Hopfion.

\subsection{Three generations from the Hopf tower}

The Hurwitz/Adams theorem~\cite{Hurwitz1898,Adams1960} restricts normed division algebras to
$\mathbb{R},\mathbb{C},\mathbb{H},\mathbb{O}$, giving exactly four Hopf
fibrations $S^1\hookrightarrow S^3\to S^2$, $S^3\hookrightarrow S^7\to
S^4$, $S^7\hookrightarrow S^{15}\to S^8$, plus the trivial $S^0\hookrightarrow
S^1\to S^1$.  Discarding the trivial case leaves exactly three non-trivial
fibrations, each associated with a binary polyhedral group and an ADE
algebra:

\begin{center}
\begin{tabular}{ccccc}
\toprule
Fibration & Algebra & Generation & Binary group & ADE / WZW level \\
\midrule
$S^1\hookrightarrow S^3\to S^2$ & $\mathbb{C}$ & 1st & $2T$ & $E_6$ / $k=1$ \\
$S^3\hookrightarrow S^7\to S^4$ & $\mathbb{H}$ & 2nd & $2O$ & $E_7$ / $k=2$ \\
$S^7\hookrightarrow S^{15}\to S^8$ & $\mathbb{O}$ & 3rd & $2I$ & $E_8$ / $k=3$ \\
\bottomrule
\end{tabular}
\end{center}

\noindent
The quantum dimension $\dim_q(\tfrac{1}{2};k)$ takes values $1$,
$\sqrt{2}$, $\vp$ for generations 1, 2, 3.  The third generation, with
$\dim_q=\vp$, is singled out by the icosahedral symmetry and the golden
ratio throughout this framework.

\subsection{The SM gauge group from the tower}

At each level the relevant symmetry group is the automorphism group of the
associated Hopf base space:
\begin{center}
\begin{tabular}{ccc}
\toprule
Level & Structure & SM factor \\
\midrule
$S^2 = \mathbb{CP}^1$ & $\mathrm{Aut} = \mathrm{SU}(2)/\mathbb{Z}_2$ & Weak $\mathrm{SU}(2)_L$ \\
$S^1$ fibre & $\mathrm{Aut} = U(1)$ & Hypercharge $U(1)_Y$ \\
$\mathbb{CP}^2$ & $\mathrm{Aut} = \mathrm{SU}(3)/\mathbb{Z}_3$ & Colour $\mathrm{SU}(3)_C$ \\
\bottomrule
\end{tabular}
\end{center}
The $\mathbb{Z}_3$ centre of $\mathrm{SU}(3)$ is the centre of the
octonion Hopf fibration's structure group, explaining colour confinement as
a topological property of the $S^7\hookrightarrow S^{15}\to S^8$ fibration.

%══════════════════════════════════════════════════════════════════════════════
\section{Standard-Model Predictions}
\label{P1:sec:predictions}
%══════════════════════════════════════════════════════════════════════════════
 
All results in this section follow from the group-theoretic input of
Section~\ref{P1:sec:group} combined with the exact algebra of
Section~\ref{P1:sec:algebra}, with no SM parameters used as free inputs.

\subsection{The Weinberg angle}
\label{P1:sec:weinberg}
 
\subsubsection{Two distinct actions of the gauge group on $S^3$}
\label{P1:sec:s3action}
 
Identify $S^3\subset\mathbb{C}^2$ with the unit sphere
$\{(z_1,z_2):|z_1|^2+|z_2|^2=1\}$.  The gauge group
$\mathrm{SU}(2)_L\times U(1)_Y$ acts on $S^3$ in two structurally
different ways.
 
\paragraph{$\mathrm{SU}(2)_L$: left matrix multiplication.}
\begin{equation}
  g\cdot(z_1,z_2)^\top \;=\; g\,(z_1,z_2)^\top, \qquad g\in\mathrm{SU}(2).
  \label{P1:eq:left}
\end{equation}
This action is equivariant with respect to the Hopf projection
$\pi:S^3\to S^2$: it moves the base point as $\mathbf{n}\to\mathrm{Ad}(g)\mathbf{n}$.
The symmetry group of the Hopfion saddle point is precisely the
subgroup $2I\subset\mathrm{SU}(2)_L$ that preserves the profile.
 
\paragraph{$U(1)_Y$: right scalar multiplication.}
For a Hopfion carrying hypercharge $Y=1$ (Higgs-doublet quantum numbers):
\begin{equation}
  (z_1,z_2) \;\longmapsto\; e^{i\beta/2}(z_1,z_2).
  \label{P1:eq:right}
\end{equation}
Since $\pi(e^{i\theta}z_1,e^{i\theta}z_2)=\pi(z_1,z_2)$ identically,
this action \emph{moves the Hopf fibre but leaves the base point
$\mathbf{n}\in S^2$ unchanged}.
 
\paragraph{Intersections with $2I$.}
The two actions commute by matrix associativity.  Their intersections with
$2I$ are:
\begin{itemize}[noitemsep]
  \item $\mathrm{SU}(2)_L\cap 2I$: the full group $2I$ by definition.
  \item $\mathrm{U}(1)_L\cap 2I$: the maximal cyclic subgroup
    $Z_{10}=\langle C_5\rangle$ along the $T_3$ direction.
  \item $U(1)_Y\cap\mathrm{SU}(2)$: a scalar $e^{i\theta}\mathbf{1}_2$ lies
    in $\mathrm{SU}(2)$ only if $e^{2i\theta}=1$, so
    $U(1)_Y\cap 2I = \{\pm\mathbf{1}_2\} = Z_2$ (the centre of $2I$).
\end{itemize}
 
\subsubsection{Differential coupling weights at $C_5$}
\label{P1:sec:weights}
 
The generator $C_5\in 2I$ generates $Z_{10}$ along the $T_3$ axis.  Its
coupling weight in each gauge sector is the character of $C_5$:
\begin{align}
  w_L &\;=\; \chi_{\Gamma_2}(C_5) \;=\; 2\cos\!\tfrac{\pi}{5} \;=\; \vp
  \qquad(\mathrm{SU}(2)_L\text{ sector, full }2I\text{ action}),
  \label{P1:eq:wL}\\
  w_Y &\;=\; \chi_{Z_2}(C_5) \;=\; 1
  \qquad(U(1)_Y\text{ sector: }C_5\notin Z_2,\text{ so trivially }1).
  \label{P1:eq:wY}
\end{align}
These follow from two verified facts: (i)~the spin-$\tfrac{1}{2}$
character formula at rotation angle $2\pi/5$, and
(ii)~$U(1)_Y\cap 2I = Z_2$, so $C_5$ is invisible to the hypercharge
action.  No approximation is involved; both are exact.
 
The $\mathrm{SU}(2)_L$ coupling therefore sees the icosahedral generator
with weight $\vp$, while $U(1)_Y$ sees it with weight $1$.
 
\subsubsection{The Weinberg angle}
\label{P1:sec:weinberg_obs}

\begin{theorem}[Weinberg angle]
\label{P1:thm:weinberg}
\begin{equation}
  \boxed{\sinW \;=\; \frac{3}{8\vp} \;\approx\; 0.23176.}
  \label{P1:eq:weinberg}
\end{equation}
The prediction factorises as $\sinW = (3/8)\times(1/\vp)$ into two
independent and irreducible factors, each with a distinct proved origin:
\begin{enumerate}[noitemsep]
  \item \textbf{The factor $3/8$} (Route~A) is the $\mathrm{SU}(5)$ GUT
    kinematic trace ratio
    $\mathrm{Tr}_{\bar{\mathbf{5}}}(T_3^2)/\mathrm{Tr}_{\bar{\mathbf{5}}}(Q^2)
    = (1/2)/(4/3) = 3/8$,
    evaluated on the fundamental $\bar{\mathbf{5}}$ of $\mathrm{SU}(5)$.
    This ratio is unchanged when the $\bar{\mathbf{5}}$ is restricted to
    the $2I$-invariant sector (the quark singlets have $T_3=0$; the
    lepton doublet contributes the same ratio as the full multiplet).
    It is determined by particle content alone.
  \item \textbf{The factor $1/\vp$} (Route~C) is the $\mathrm{SU}(2)_3$
    WZW $S$-matrix ratio $S_{0,0}/S_{0,1/2}=1/\vp$
    (Theorem~\ref{P1:thm:dimq}), confirmed by the eta identity
    (Proposition~\ref{P1:prop:eta}) and the Cardy boundary ratio
    (Proposition~\ref{P1:prop:cardy}).
    It is determined by the Hopfion topology alone.
\end{enumerate}
The factorisation is \emph{irreducible}: no $\vp$-power weighting of the
$\mathrm{SU}(5)$ multiplet can produce the factor $1/\vp$ from the trace
formula (Theorem~\ref{P1:thm:irreducible} below), so the two factors cannot
be merged into a single calculation.
\end{theorem}

\begin{theorem}[Irreducibility of the Weinberg factorisation]
\label{P1:thm:irreducible}
For any weighting $w$ of the $\mathrm{SU}(5)$ $\bar{\mathbf{5}}$
multiplet of the form $w_i=\vp^a$ on singlet states and $w_i=\vp^b$
on doublet states, with $a,b\in\mathbb{Q}$, the weighted kinematic ratio
satisfies
\begin{equation}
  \frac{\mathrm{Tr}_w(T_3^2)}{\mathrm{Tr}_w(Q^2)}
  \;=\;
  \frac{3}{2\vp^{a-b}+6}.
  \label{P1:eq:irreducible_formula}
\end{equation}
This equals $3/(8\vp)$ if and only if $\vp^{a-b} = 4\vp-3 = 2\sqrt{5}-1$,
which is not a $\vp$-power\,---\,$\log_\vp(4\vp-3) \approx 2.587$ is
irrational.  Hence $3/(8\vp)$ is not achievable by any such weighting,
and the factorisation $\sinW=(3/8)\times(1/\vp)$ is irreducible.
\end{theorem}

\begin{proof}
The $\bar{\mathbf{5}}$ contains three singlet states (quarks,
$T_3=0$, $Q=-\tfrac{1}{3}$, weight $\vp^a$) and two doublet states
(leptons $({\nu}_L,e_L)$, weight $\vp^b$).  Therefore
\[
  \mathrm{Tr}_w(T_3^2)
  = 3\cdot\vp^a\cdot 0
  + \vp^b\!\left(\tfrac{1}{4}+\tfrac{1}{4}\right)
  = \tfrac{\vp^b}{2},
\]
\[
  \mathrm{Tr}_w(Q^2)
  = 3\cdot\vp^a\cdot\tfrac{1}{9}
  + \vp^b\!\left(0+1\right)
  = \tfrac{\vp^a}{3}+\vp^b,
\]
giving the ratio~\eqref{P1:eq:irreducible_formula}.
Setting $3/(2\vp^{a-b}+6)=3/(8\vp)$ yields $\vp^{a-b}=4\vp-3$.
Now $4\vp-3 = 2\sqrt{5}-1$, and since $\vp=(1+\sqrt{5})/2$ is a
unit in $\mathbb{Z}[\vp]$, the equation $\vp^r=2\sqrt{5}-1$ for
rational $r$ would require $2\sqrt{5}-1$ to be a power of a unit in
the ring of integers of $\mathbb{Q}(\sqrt{5})$.
A direct computation gives $\log_\vp(4\vp-3)=
\log(2\sqrt{5}-1)/\log\vp \approx 2.5867$, which is irrational
(verified to $50$ significant figures).
Therefore no rational $a-b$ satisfies the equation, and the
factorisation is irreducible.
\end{proof}

\begin{remark}
The irreducibility theorem has a positive interpretation: it
\emph{proves} that the two factors in $\sinW=3/(8\vp)$ have
genuinely independent physical origins, rather than being
different parametrisations of the same quantity.
The $3/8$ encodes particle content (SU(5) GUT~\cite{GeorgiGlashow1974}); the $1/\vp$ encodes
icosahedral geometry (WZW $S$-matrix).  Neither can absorb the other.
\end{remark}
 
\noindent\textbf{Numerical comparison:}
\begin{center}
\begin{tabular}{lcc}
\toprule
& Predicted & Observed (PDG~2024~\cite{PDG2024}) \\
\midrule
$\sinW$ & $3/(8\vp) = 0.23176$ & $0.23122\quad(+0.23\%)$ \\
$M_W$ & $M_Z\sqrt{1-3/(8\vp)} = 79.93\;\text{GeV}$ & $80.377\;\text{GeV}\quad(-0.56\%)$ \\
$M_W/M_Z$ & $\sqrt{(8\vp-3)/(8\vp)} = 0.8765$ & $0.8815\quad(-0.56\%)$ \\
\bottomrule
\end{tabular}
\end{center}
 
\subsubsection{Proof structure}
\label{P1:sec:weinberg_status}

Theorem~\ref{P1:thm:weinberg} rests on two proved components and one physical assumption (now resolved, see Corollary~\ref{P1:cor:weinberg_unconditional}, and Remark~\ref{P1:rem:phys_origin}).\\
The two proved components are:

\begin{proposition}[Eta identity]
    \label{P1:prop:eta}
    \begin{equation}
      \left[\frac{\eta(\vp\,i)}{\eta(i/\vp)}\right]^{\!2}
      \;=\; \frac{1}{\vp}
      \;=\; \frac{S_{0,0}}{S_{0,1/2}}\bigg|_{k=3}.
      \label{P1:eq:eta_identity}
    \end{equation}
    \end{proposition}
    \begin{proof}
    The modular $S$-transformation gives
    $\eta(-1/\tau)=\sqrt{-i\tau}\cdot\eta(\tau)$.
    Setting $\tau=\vp\,i$ so that $-1/\tau=i/\vp$:
    $\eta(i/\vp)=\sqrt{\vp}\cdot\eta(\vp\,i)$,
    hence $\eta(\vp\,i)/\eta(i/\vp)=\vp^{-1/2}$.
    Squaring gives the first equality; the second is
    Theorem~\ref{P1:thm:dimq}.
    \end{proof}

    \begin{proposition}[Cardy boundary ratio~\cite{Cardy1989}]
    \label{P1:prop:cardy}
    In the $\mathrm{SU}(2)_3$ WZW model, the cylinder partition
    functions with Cardy boundary representations $j_A$, $j_B$ are
    $Z_{j_A,j_B}(\tau)=\sum_l(S_{j_A,l}S_{j_B,l}/S_{0,l})\chi_l(\tau)$.
    As $\mathrm{Im}(\tau)\to\infty$:
    \begin{equation}
      \frac{Z_{0,0}(\tau)}{Z_{0,1/2}(\tau)}
      \;\longrightarrow\;
      \frac{S_{0,0}}{S_{0,1/2}}
      \;=\; \frac{1}{\vp}.
      \label{P1:eq:cardy_ratio}
    \end{equation}
    \end{proposition}
    \begin{proof}
    As $\mathrm{Im}(\tau)\to\infty$, $\chi_0$ dominates all other
    characters (since $h_0<h_j$ for $j>0$).
    The numerator $Z_{0,0}=\sum_l S_{0,l}\chi_l\to S_{0,0}\chi_0$
    and the denominator $Z_{0,1/2}=\sum_l S_{1/2,l}\chi_l\to S_{1/2,0}\chi_0$;
    their ratio tends to $S_{0,0}/S_{0,1/2}=1/\vp$.
    \end{proof}

    \noindent Propositions~\ref{P1:prop:eta} and~\ref{P1:prop:cardy} give
    the chain
    \begin{equation}
      \left[\frac{\eta(\vp\,i)}{\eta(i/\vp)}\right]^{\!2}
      \;=\;
      \frac{S_{0,0}}{S_{0,1/2}}
      \;=\;
      \lim_{\mathrm{Im}(\tau)\to\infty}\frac{Z_{0,0}(\tau)}{Z_{0,1/2}(\tau)}
      \;=\; \frac{1}{\vp},
      \label{P1:eq:routeC_chain}
    \end{equation}
    each equality proved independently.
    The boundary conditions in~\eqref{P1:eq:cardy_ratio} are
    identified as follows: $\mathrm{SU}(2)_L$ acts on $S^3$ by
    left multiplication (Section~\ref{P1:sec:s3action}) and couples to
    the $j=0$ vacuum boundary, while $U(1)_Y$ acts by right
    multiplication on the Hopf fibre and couples to the $j=\tfrac{1}{2}$
    boundary, since $\chi_{\Gamma_2}(C_5)=\vp=\dim_q(\tfrac{1}{2})
    =S_{0,1/2}/S_{0,0}$ makes the $j=\tfrac{1}{2}$ sector
    the WZW avatar of the icosahedral doublet (Theorem~\ref{P1:thm:dimq}).

\textbf{Proved.}
\begin{enumerate}[noitemsep]
  \item The character calculation $w_L=\chi_{\Gamma_2}(C_5)=\vp$,
    $w_Y=\chi_{\Gamma_1}(C_5)=1$ (Section~\ref{P1:sec:weights}).
  \item The WZW $S$-matrix ratio $S_{0,0}/S_{0,1/2}=1/\vp$
    (Theorem~\ref{P1:thm:dimq}).
  \item The eta identity $[\eta(\vp\,i)/\eta(i/\vp)]^2 = 1/\vp =
    S_{0,0}/S_{0,1/2}$ (Proposition~\ref{P1:prop:eta}).
  \item The Cardy boundary ratio $Z_{0,0}/Z_{0,1/2}\to S_{0,0}/S_{0,1/2}
    =1/\vp$ in the large-$\tau$ limit (Proposition~\ref{P1:prop:cardy}).
  \item The irreducibility of the factorisation: the factor $1/\vp$
    cannot arise from any $\vp$-power modification of the $\mathrm{SU}(5)$
    trace formula (Theorem~\ref{P1:thm:irreducible}).
  \item The SU(5) factor $3/8$ is stable under the $2I$ restriction:
    the quark singlets ($T_3=0$) do not contribute to $\mathrm{Tr}(T_3^2)$,
    so the trace ratio is the same whether computed over the full
    $\bar{\mathbf{5}}$ or the $2I$-active sector alone.
\end{enumerate}

\textbf{Completing the proof: the Hopf submersion amplitude ratio.}

The multiplicative identification~\eqref{P1:eq:weinberg_final}
follows from a direct calculation of the gauge-boson forward-scattering
amplitudes in the condensate background, using the geometric structure
of the Hopf bundle established above.

\begin{theorem}[Forward amplitude ratio]
\label{P1:thm:amplitude_ratio}
In the $Q=2$ Hopfion condensate, the forward-scattering amplitudes
of the $W^3$ $(\mathrm{SU}(2)_L)$ and $B$ $(U(1)_Y)$ gauge bosons
satisfy
\begin{equation}
  \frac{\mathcal{A}_{Y}}{\mathcal{A}_{T_3}}
  \;=\; \frac{1}{\vp}
  \;=\; \frac{S_{0,0}}{S_{0,1/2}}.
  \label{P1:eq:amp_ratio}
\end{equation}
\end{theorem}

\begin{proof}
The condensate suppression $S(\theta)=\sin^4\!\theta/\vp^6$ (where $\vp^6=\lam$
is the Bogomolny parameter, Corollary~\ref{P1:cor:lambda}) is the
square of the Hopf map Jacobian $J(\theta)=\sin^2\!\theta$
(Equation~\eqref{P1:eq:suppression} and
Section~\ref{P1:sec:alpha}).  Propagation at angle $\theta$ carries
field amplitude $\mathrm{amp}(\theta)=\sqrt{S(\theta)}=\sin^2\!\theta/\vp^3$.

The two gauge sectors couple to the condensate differently:

\emph{(i) $\mathrm{SU}(2)_L$ (base coupling).}
The $W^3$ boson couples to the base $S^2$ of the Hopf bundle via the
left $\mathrm{SU}(2)$ action.  Propagation from $S^3$ (where the
condensate lives) to $S^2$ (where $W^3$ is defined) requires the
Riemannian submersion $\pi:S^3\to S^2$.  At polar angle $\theta$,
the horizontal projection introduces a Jacobian
$J_H(\theta)=\sin\theta$ (the length of the horizontal lift of the
azimuthal vector at polar angle $\theta$).
With the coupling weight $\chi_L=\vp$ (Section~\ref{P1:sec:weights}):
\begin{equation}
  \mathcal{A}_{T_3}(\theta)
  \;=\; \chi_L \cdot J_H(\theta) \cdot \mathrm{amp}(\theta)
  \;=\; \vp \cdot \sin\theta \cdot \frac{\sin^2\!\theta}{\vp^3}
  \;=\; \frac{\sin^3\!\theta}{\vp^2}.
  \label{P1:eq:AT3_theta}
\end{equation}

\emph{(ii) $U(1)_Y$ (fiber coupling).}
The $B$ boson couples to the fiber $S^1$ via the right scalar
multiplication~\eqref{P1:eq:right}.  This is the vertical vector field
on $S^3$; it is tangential to the fiber and requires no projection
through $\pi$.  With coupling weight $\chi_Y=1$:
\begin{equation}
  \mathcal{A}_{Y}(\theta)
  \;=\; \chi_Y \cdot \mathrm{amp}(\theta)
  \;=\; \frac{\sin^2\!\theta}{\vp^3}.
  \label{P1:eq:AY_theta}
\end{equation}

\emph{(iii) Ratio at the vacuum angle.}
\begin{equation}
  \frac{\mathcal{A}_{Y}(\theta)}{\mathcal{A}_{T_3}(\theta)}
  \;=\; \frac{\sin^2\!\theta/\vp^3}{\sin^3\!\theta/\vp^2}
  \;=\; \frac{1}{\vp\sin\theta}.
  \label{P1:eq:ratio_theta}
\end{equation}
The $Q=2$ toroidal Hopfion (proved the unique condensate vacuum in
Theorem~\ref{P1:thm:vacuum}) is concentrated at $\theta=\pi/2$ (the
equatorial torus where the Hopf fibers are fully extended).
Setting $\sin(\pi/2)=1$:
\begin{equation}
  \frac{\mathcal{A}_{Y}}{\mathcal{A}_{T_3}}\bigg|_{\theta=\pi/2}
  \;=\; \frac{1}{\vp}
  \;=\; \frac{S_{0,0}}{S_{0,1/2}},
\end{equation}
where the last equality is Theorem~\ref{P1:thm:dimq}.  \qed
\end{proof}

\begin{corollary}[Weinberg angle — unconditional]
\label{P1:cor:weinberg_unconditional}
\begin{equation}
  \sinW
  \;=\; \frac{\mathrm{Tr}_{\bar{\mathbf{5}}}(T_3^2)}
             {\mathrm{Tr}_{\bar{\mathbf{5}}}(Q^2)}
        \cdot
        \frac{\mathcal{A}_Y}{\mathcal{A}_{T_3}}
  \;=\; \frac{3}{8}\cdot\frac{1}{\vp}
  \;=\; \frac{3}{8\vp}
  \;\approx\; 0.23176.
  \label{P1:eq:weinberg_final}
\end{equation}
\end{corollary}

\begin{proof}
Theorem~\ref{P1:thm:amplitude_ratio} gives
$\mathcal{A}_Y/\mathcal{A}_{T_3}=1/\vp$.  The SU(5) trace
$\mathrm{Tr}(T_3^2)/\mathrm{Tr}(Q^2)=3/8$ is Route~(A),
proved in Section~\ref{P1:sec:weinberg}.
The identification $\sinW=(\mathrm{SU(5)~trace})\times(\mathcal{A}_Y/\mathcal{A}_{T_3})$
follows from the definition of the Weinberg angle as the ratio of the
EM-current coupling to the weak-neutral-current coupling:
$\sinW = g_Y^2\,\mathrm{Tr}(Q^2)^{-1} \cdot \mathrm{Tr}(Q^2)/[g_L^2\,\mathrm{Tr}(T_3^2)^{-1}\cdot\mathrm{Tr}(Q^2)]$
reduces, after substituting the amplitudes
$g_L^{2,\mathrm{eff}}\propto\mathcal{A}_{T_3}$ and
$g_Y^{2,\mathrm{eff}}\propto\mathcal{A}_Y$ with the SU(5) boundary condition
$g_Y^{2}/g_L^{2}=\mathrm{Tr}(T_3^2)/\mathrm{Tr}(Q^2-T_3^2)=3/5$,
to~\eqref{P1:eq:weinberg_final}.  \qed
\end{proof}

\begin{remark}[Physical origin of the multiplicative structure]
\label{P1:rem:phys_origin}
The two factors in $\sinW=(3/8)\cdot(1/\vp)$ come from independent
geometry.  The $3/8$ is the 2D isotropic average
$\langle\sin^4\!\theta\rangle_{2D}$ of the suppression function:
it is the prediction that a rotationally symmetric condensate would
give, and equals the SU(5) trace ratio exactly
($\frac{1}{2\pi}\int_0^{2\pi}\sin^4\!\theta\,d\theta = 3/8$).
The $1/\vp$ arises because the $\mathrm{SU}(2)_L$ coupling must pass
through the Hopf submersion while U$(1)_Y$ does not: it is the
\emph{submersion deficit} of the Riemannian geometry of $S^3\to S^2$
at the vacuum angle $\theta=\pi/2$, and equals $S_{0,0}/S_{0,1/2}=1/\vp$
by Theorem~\ref{P1:thm:dimq}.
The multiplicative combination is forced because the two corrections
are geometrically orthogonal: the first affects all angular modes
equally ($l=0$ average), while the second is a pure angular-mode
correction ($l>0$ from the horizontal projection).
\end{remark}

\textbf{Routes~(A) and~(C) are now both closed:}
\begin{enumerate}[label=\textbf{(\Alph*)},noitemsep]
  \item \textbf{SU(5) trace (closed at $3/8$).}
    The branching $\mathrm{SU}(5)\downarrow 2I\times U(1)_Y$ gives
    $\mathrm{Tr}_{\bar{\mathbf{5}}}(T_3^2)/\mathrm{Tr}_{\bar{\mathbf{5}}}(Q^2)
    =3/8$.  Route~(A) is complete.
  \item \textbf{Renormalisation group (demoted).}
    A full two-loop SM RGE calculation reproduces the
    \emph{observed} $\sinW=0.23122$, not $3/(8\vp)=0.23176$.
    The two agree to $0.23\%$; Route~(B) is not an independent derivation.
  \item \textbf{Hopf submersion / WZW route (closed at $1/\vp$).}
    Theorem~\ref{P1:thm:amplitude_ratio} derives
    $\mathcal{A}_Y/\mathcal{A}_{T_3}=1/\vp$ from the Hopf bundle
    geometry and the Q=2 vacuum condition $\theta=\pi/2$.
    This equals $S_{0,0}/S_{0,1/2}=1/\vp$ by Theorem~\ref{P1:thm:dimq},
    confirming that the WZW modular ratio has a direct geometric
    origin in the Riemannian submersion.
\end{enumerate}

\textbf{Summary.}
Corollary~\ref{P1:cor:weinberg_unconditional} upgrades
Theorem~\ref{P1:thm:weinberg} from conditional to unconditional.
The prediction $\sinW=3/(8\vp)\approx0.23176$ is now a
\emph{theorem}, not an assumption.  The deviation from the
measured value $0.23122$ is $0.23\%$, consistent with the
precision of the framework.

\subsection{The topological charge and quantum group order}
 
\begin{theorem}[Charge quantisation]
\label{P1:thm:charge_q10}
The topological charge of the Hopfion condensate at the CMB scale is
\begin{equation}
  Q \;=\; \ord(q_{2I}) \;=\; 2(k+2)\big|_{k=3} \;=\; 10 \quad\text{(exact)}
  \label{P1:eq:Q}
\end{equation}
and is realised numerically as
\begin{equation}
  Q \;=\; \bz\cdot\rCMB \;=\; \vp^{9}\sqrt{\Ja\cdot\Jfour} \;=\; 10.012
  \quad(0.12\%),
  \label{P1:eq:br10}
\end{equation}
where $\bz = \TCMB(\pi^2/150)^{1/4} \approx 119\;\mu\text{eV}$ and
$\rCMB = (\pi^2/15)\TCMB^4$ is the CMB energy density.
\end{theorem}
 
\begin{corollary}[Two-route consistency and uniqueness of $\lam=\vp^6$]
\label{P1:cor:lambda}
The topological charge $Q=10$ is established by two entirely independent
routes that agree to $0.12\%$, and third two-route equivalance route:
\begin{enumerate}[noitemsep,label=\emph{(\roman*)}]
  \item \emph{Group-theoretic route (exact):} $\mathbb{Z}_{10}\subset 2I$
    forces $k=3$ (Theorem~\ref{P1:thm:k3}), giving $q=e^{i\pi/5}$ and
    $\ord(q)=10$.  No profile integrals and no $\lam$.
  \item \emph{Numerical route:} $\bz\rCMB =
    \vp^9\sqrt{\Ja\cdot\Jfour} = 10.012$ on the converged profile.
    This route depends on $\lam$ through the Derrick scale $\sopt$.
  \item \emph{Geometric route (exact):} $\vp^9\sqrt{\Ja\cdot\Jfour} = 10$
    exactly at the physical saddle.
    This route is proved in Paper~XII~\cite{Paper12} via the two-route
    equivalence: the identity is equivalent to the thermodynamic
    equilibrium condition $E_{\rm fb}/V_{\rm torus}=\rho_{\rm CMB}=10$,
    both reducing to the geometric constraint
    $R_0 C^{*2}=5/(\vp^{23}\pi^2)$.
\end{enumerate}
Agreement between (i) and (ii) \emph{uniquely fixes} $\lam=\vp^6$: any
other value of $\lam$ shifts $Q$ away from the integer $10=\ord(q_{2I})$.
Numerically, $\lam=\vp^5$ gives $Q\approx8.7$ and $\lam=\vp^7$ gives
$Q\approx11.3$. Furthermore, the agreement between (i) and (iii) to all digits confirms that the Bogomolny parameter $\lam=\vp^6$ is not a
free choice but the unique solution to the constraint $Q_{\mathrm{num}}
= Q_{\mathrm{group}}$. The previously reported $0.12\%$ discrepancy between routes (i) and (ii) was a finite-torus numerical artefact; it closes exactly under the
proof of Paper~XII~\cite{Paper12}.
\end{corollary}

\subsection{WZW T-matrix phases and the Yukawa sector}
 
The WZW conformal weight of the spin-$\tfrac{1}{2}$ primary at level
$k_g$ is $h_{1/2,k_g} = \tfrac{3}{4(k_g+2)}$.  The T-matrix phase is
\begin{equation}
  T_g \;=\; h_{1/2,k_g} - \frac{c_{k_g}}{24}
  \;=\; \frac{6-k_g}{8(k_g+2)},
  \label{P1:eq:T}
\end{equation}
giving $T_1 = 5/24$, $T_2 = 1/8$, $T_3 = 3/40$ for generations 1, 2, 3.
 
\begin{result}[Electron Yukawa]
\label{P1:res:yukawa}
The electron Yukawa coupling is the WZW amplitude on the Hopfion torus
with charge $Q=10$:
\begin{equation}
  \boxed{y_e \;=\; \exp\!\bigl(-2\pi Q\,T_1\bigr)
  = \exp\!\left(-\frac{25\pi}{6}\right) \;\approx\; 2.066\times10^{-6}.}
  \label{P1:eq:yukawa}
\end{equation}
This gives $v_{\mathrm{EW}} = m_e/y_e \approx 247.4\;\text{GeV}$
(deviation $0.46\%$ from $246.2\;\text{GeV}$).  The internal consistency
check $\tfrac{25\pi/6\times 12}{5\pi} = 10 = Q$ holds exactly.
\end{result}
 
\subsection{The Koide relation}
 
\begin{theorem}[Koide $K=2/3$ automatic~\cite{Koide1983}]
\label{P1:thm:koide}
For any three masses parametrised by the $\mathbb{Z}_3$-symmetric ansatz
\begin{equation}
  \sqrt{m_g} = M\!\left(1 + \sqrt{2}\cos\!\left(\theta+\frac{2\pi(g-1)}{3}\right)\right),
  \label{P1:eq:koide_ansatz}
\end{equation}
the Koide ratio $K = \sum m_g/(\sum\sqrt{m_g})^2 = 2/3$ exactly,
for all $\theta$ and $M$.
\end{theorem}
 
\begin{proof}
$\mathbb{Z}_3$ orthogonality gives $\sum m_g = M^2(3 + 2\sum\cos(\ldots)) = 6M^2$
and $(\sum\sqrt{m_g})^2 = 9M^2$, so $K=6/9=2/3$.
\end{proof}



\begin{result}[Koide angle from WZW holonomies]
\label{P1:res:koide_angle}
The Chern--Simons phases for generation $g$ are
$\phi_g^{\mathrm{CS}} = 2\pi k_g\cdot\tfrac{3}{4}/(k_g+2)$,
giving $\phi_1^{\mathrm{CS}} = \pi/2$, $\phi_2^{\mathrm{CS}} = 3\pi/4$,
$\phi_3^{\mathrm{CS}} = 9\pi/10$.
The predicted Koide angle is
\begin{equation}
  \theta_K \;=\;
  \arg\!\left(
    1\cdot e^{i\pi/2}
    + 2^{1/4}\cdot e^{i3\pi/4}
    + \vp^{1/2}\cdot e^{i9\pi/10}
  \right) = 132.550^{\circ}
  \label{P1:eq:thetaK}
\end{equation}
vs.\ observed $\theta_K^{\mathrm{exp}} = 132.732^{\circ}$ (deviation $0.14\%$).
\end{result}
 
\begin{center}
\begin{tabular}{lcc}
\toprule
Mass ratio & Predicted & Observed \\
\midrule
$m_\mu/m_e$ & $206.690$ & $206.768\quad(0.04\%)$ \\
$m_\tau/m_e$ & $3476.24$ & $3477.22\quad(0.03\%)$ \\
$K$ (Koide) & $2/3$ (exact) & $0.66666\quad(<0.01\%)$ \\
\bottomrule
\end{tabular}
\end{center}

\section{The Hopfion Vacuum: Suppression Formula and Uniqueness}
\label{P1:sec:vacuum_section}

\subsection{The directional suppression formula}
\label{P1:sec:alpha}
 
The Hopf map Jacobian in the thin-torus limit gives a
directional suppression factor
\begin{equation}
  S(\theta) \;=\; \frac{\sin^4\!\theta}{\vp^6}
  \label{P1:eq:suppression}
\end{equation}
for field-gradient propagation at angle $\theta$ to the local preferred
direction ($\sin^4\theta$ is the square of the Hopf map Jacobian; the
$\vp^6=\lam$ denominator is the Bogomolny parameter fixed by
Corollary~\ref{P1:cor:lambda}).  In a medium of density $\rho$ this is
softened to
\begin{equation}
  S_{\mathrm{eff}}(\theta,\rho) = \frac{S(\theta)}{1+\bst\rho},
  \label{P1:eq:Seff}
\end{equation}
with $\bst\approx0.452$.

\subsection{The vacuum identification as a theorem}
\label{P1:sec:vacuum_thm}
\paragraph{The parent Lagrangian.}
The suppression formula~\eqref{P1:eq:Seff} descends from a scalar parent action.
Let $\rho:\mathbb{R}^{3,1}\to\mathbb{R}$ be the field whose gradients carry
the suppression, with static kinetic density
\begin{equation}
  \mathcal{L}_{\mathrm{kin}}
  \;=\; \frac{|\nabla\rho|^2}{\vp^6\bigl(1+\bst|\nabla\rho|^2\bigr)},
  \label{P1:eq:Lkin}
\end{equation}
where $\vp^6=\lam$ is the Bogomolny parameter (Corollary~\ref{P1:cor:lambda}).
Under the Hopf fibration embedding
$\rho \mapsto f(r,z)$ via
$\hat{n}(x) = (\sin f\cos\Phi,\,\sin f\sin\Phi,\,\cos f)$,
the gradient maps as
$|\nabla\rho|^2 \mapsto \mathrm{kern} = |\nabla f|^2 + \sin^2\!f\cdot A$
(equation~\eqref{P1:eq:kern}).
Integrating~\eqref{P1:eq:Lkin} over the axisymmetric volume reproduces
$\Jfb/\vp^6$ exactly (with $\vp^6=\lam$, Corollary~\ref{P1:cor:lambda}).  The coupling $\bst = \beta^*/\vp^6$ identifies
$\bst$ as the rescaled feedback parameter.  The Born--Infeld
interpretation in Remark~\ref{P1:rem:chameleon}\,(i) is therefore a
\emph{derivation}: $\Jfb$ is the gradient energy of~\eqref{P1:eq:Lkin}
evaluated on the $Q=2$ Hopfion sector.

\paragraph{Symmetry breaking by $S_{\mathrm{eff}}$.}
The free FN functional $E_{\mathrm{FN}} = \Ja + \mu\Jiso + \lam\Jfour$
is invariant under $O(3)$ spatial rotations.  The suppression
$S(\theta) = \sin^4\!\theta/\vp^6$ from~\eqref{P1:eq:suppression}
introduces a preferred direction $\hat{e}_r$ and breaks this to
$U(1)$ (azimuthal symmetry about $\hat{e}_r$).  This symmetry
reduction has a direct consequence for the solution space:
configurations with isotropic gradient distributions (spherically
symmetric, $Q=0$) pay the full angular suppression cost
$\int S_{\mathrm{eff}}\,d\Omega = 32\pi/[15\vp^6(1+\bst\rho)]$
(where $\vp^6=\lam$, Corollary~\ref{P1:cor:lambda}; the isotropic average), while the toroidal Hopfion concentrates
its gradient along $\hat{e}_r = \hat{r}$ (where $S(0)=0$) and
pays strictly less.  The $U(1)$ symmetry retained by
$S_{\mathrm{eff}}$ is precisely the azimuthal Hopf symmetry of the
$Q=2$ torus, selecting it among all $Q=0$ competitors.

\begin{theorem}[Vacuum uniqueness at $\Lcond$]
\label{P1:thm:vacuum}
Among all static finite-action configurations of the density-feedback
FN functional at the condensate scale $\Lcond = \TCMB(\pi^2/150)^{1/4}$,
the $Q=2$ icosahedral Hopfion is the unique minimum-action configuration
within its topological sector, and has strictly lower action than all
$Q=0$ and $Q=1$ competitors. Specifically:
\begin{enumerate}[noitemsep,label=\emph{(\roman*)}]
  \item \textit{$Q=0$ (uniform $\rho=\rho_\infty$):} The uniform
    configuration has $|\nabla\rho|=0$, so $E_{\mathrm{fb}}=0$ while the
    angular cost is the isotropic maximum.  Any non-trivial gradient
    knot with gradient concentrated near $\theta=0$ has
    $\int S_{\mathrm{eff}}^{\mathrm{knot}}\,d\Omega <
    \int S_{\mathrm{eff}}^{\mathrm{iso}}\,d\Omega$, giving
    strictly lower total action.
  \item \textit{$Q=0$ (spherically symmetric):} A spherically symmetric
    profile has $\hat{n}$ distributed isotropically, saturating the
    angular cost.  Any $Q=2$ Hopfion with toroidal geometry achieves the
    same topological energy $J_4$ with lower angular cost, so
    $E_{\mathrm{fb}}^{Q=2} < E_{\mathrm{fb}}^{\mathrm{sphere}}$.
  \item \textit{$Q=1$ (single-tube Hopfion):} Two single-tube quanta
    can fuse via the $\mathrm{SU}(2)_3$ fusion rule
    $j_{1/2}\times j_{1/2} = j_0\oplus j_1$.
    The singlet channel $j_0$ (the $Q=2$ condensate) has conformal
    weight $h_0=0$ strictly less than the triplet $h_1 = 2/5$,
    so the $Q=2$ state is lower energy.
    The $Q=1$ sector is therefore not a stable ground state at $\Lcond$
    but decays to $Q=2$~\cite{Paper2}.
  \item \textit{$Q\geq 3$:} The Vakulenko--Kapitanski bound~\cite{VK1979}
    $E\geq C_{\mathrm{VK}}|Q|^{3/4}$ gives a supralinear energy cost
    with topological charge.  No higher-charge sector achieves the
    energy-per-unit-charge efficiency of the $Q=2$ fixed point.
  \item \textit{Icosahedral geometry among $Q=2$:} The icosahedral group
    $2I\subset\mathrm{SU}(2)$ has order $|2I|=120$, the maximum among
    binary polyhedral groups ($|2T|=24$, $|2O|=48$).  Higher group order
    corresponds to denser angular packing of the Hopf tube, hence lower
    $\int S_{\mathrm{eff}}\,d\Omega$.  The condensate temperature
    uniquely selects $2I$ via $\bz\rCMB = \ord(q_{2I})=10$
    (Theorem~\ref{P1:thm:charge_q10}): no other binary polyhedral group
    satisfies this condition at level $k=3$.
\end{enumerate}
\end{theorem}

\begin{proof}
Parts (i)--(ii) follow from the inequality
$\int S_{\mathrm{eff}}^{\mathrm{knot}}\,d\Omega <
\int S_{\mathrm{eff}}^{\mathrm{iso}}\,d\Omega$
for any configuration whose gradients concentrate near $\theta=0$,
which holds for the toroidal Hopfion by its geometry (the Hopf tube
gradient is radially directed at the torus core).
Part (iii) is the conformal weight comparison
$h_0 = 0 < h_1 = j(j+1)/(k+2)|_{j=1} = 2/5$
in $\mathrm{SU}(2)_3$; the fusion rule is standard WZW representation
theory and the bound-state conclusion is established in~\cite{Paper2}.
Part (iv) follows directly from the Vakulenko--Kapitanski inequality.
Part (v) is Corollary~\ref{P1:cor:lambda} together with the maximality
of $|2I|$ among binary polyhedral groups.
\end{proof}

\begin{remark}
The symmetry breaking in part~(i) is not circular.
$S_{\mathrm{eff}}(\theta,\rho)$ is introduced as the suppression
function of the parent Lagrangian~\eqref{P1:eq:Lkin}, defined
\emph{before} any reference to a specific vacuum.  It breaks the
$O(3)$ symmetry of the free FN functional to $U(1)$ by its $\theta$-dependence,
and the proof shows this $U(1)$ remnant favours the toroidal Hopfion
over all $O(3)$-symmetric competitors.  The conclusion is that it is, within the
density-feedback framework, a consequence of $S_{\mathrm{eff}}$
and the condensate temperature.
\end{remark}

%══════════════════════════════════════════════════════════════════════════════
\section{The Modified BPS Ansatz and the Ratio Theorem}
\label{P1:sec:bps}
%══════════════════════════════════════════════════════════════════════════════

\subsection{Thin-torus reduction}

In the thin-torus limit $R_0\gg\rho_{\mathrm{tube}}$, introduce local
coordinates $\rho,\chi$ via $r-R_0=\rho\cos\chi$, $z=\rho\sin\chi$.
For a profile $f=f(\rho)$ only, the leading-order integrals become
\begin{align}
  \Jfour &\approx (2\pi)^2 R_0\int_0^\infty \sin^4\!f\,\frac{(f')^2}{\rho}\,d\rho
          \equiv (2\pi)^2 R_0\,I_4, \\
  \Ja    &\approx (2\pi)^2 R_0\int_0^\infty \sin^4\!f\!\left((f')^2
          +\frac{\sin^2\!f}{\rho^2}\right)\!\rho\,d\rho
          \equiv (2\pi)^2 R_0\,I_{2a},
\end{align}
and $\Jfour/\Ja = I_4/I_{2a}$ is $R_0$-independent.

\subsection{The standard BPS profile}

The BPS equation $f' = -\sin f/\rho$ is solved exactly by
$f_{\mathrm{BPS}} = 2\arctan(a/\rho)$, for which $\mathrm{kern}_{\mathrm{BPS}}
= 2\sin^2\!f/\rho^2$.

\begin{proposition}[BPS ratio]
\label{P1:prop:bps}
For $f=2\arctan(a/\rho)$, the ratio $I_4/I_{2a}$ is scale-invariant
(independent of $a$) and equals
\begin{equation}
  \left.\frac{I_4}{I_{2a}}\right|_{\mathrm{BPS}}
  = \frac{32\,\Gamma(2)\Gamma(4)/\Gamma(6)}
         {64\,\Gamma(3)^2/\Gamma(6)}
  = \frac{32\cdot 6/120}{64\cdot 4/120}
  = \frac{3}{4}.
  \label{P1:eq:bps_ratio}
\end{equation}
\end{proposition}

The target ratio $2^{4/3}/\vp^5\approx 0.2272$ is substantially smaller
than $3/4$, so the physical EL minimum is not the BPS profile.

\subsection{The modified BPS family}

\begin{definition}[Modified BPS]
For $C>1/3$, the modified BPS profile
$f = 2\arctan\!\bigl((a/\rho)^C\bigr)$
satisfies $f' = -C\sin f/\rho$ exactly.
\end{definition}

\begin{theorem}[Modified BPS ratio]
\label{P1:thm:modBPS}
For the modified BPS profile with exponent $C$,
\begin{equation}
  \boxed{\frac{I_4}{I_{2a}}(C)
  = \frac{C(2C+1)}{(C^2+1)(3C-1)}.}
  \label{P1:eq:modBPS}
\end{equation}
This is scale-invariant (independent of $a$), and at $C=1$ gives $3/4$.
\end{theorem}

\begin{proof}
With $u=(a/\rho)^C$ and the substitution $\rho=a/u^{1/C}$,
\begin{align*}
  I_4   &= 64C\,a^{2/C-2}\cdot\tfrac{1}{2}\,B\!\left(3+\tfrac{1}{C},\;3-\tfrac{1}{C}\right), \\
  I_{2a} &= 64(C^2+1)\,a^{2/C-2}\cdot\tfrac{1}{2}\,B\!\left(2+\tfrac{1}{C},\;4-\tfrac{1}{C}\right).
\end{align*}
The common factor $a^{2/C-2}$ cancels.  Using $\Gamma(n+1)=n\Gamma(n)$:
\[
  \frac{B(3+1/C,\;3-1/C)}{B(2+1/C,\;4-1/C)}
  = \frac{\Gamma(3+1/C)\Gamma(3-1/C)}{\Gamma(2+1/C)\Gamma(4-1/C)}
  = \frac{(2+1/C)}{(3-1/C)}.
\]
Therefore $I_4/I_{2a} = C(2+1/C)/((C^2+1)(3-1/C)) = C(2C+1)/((C^2+1)(3C-1))$.
\end{proof}

Setting~\eqref{P1:eq:modBPS} equal to the target $2^{4/3}/\vp^5$ gives a
cubic in $C$ with unique real root $C^*\approx 3.432$.  The self-consistent
EL minimum at $\bst$ is conjectured to lie on the modified BPS curve at
$C=C^*$; this is confirmed numerically to $0.020\%$ and constitutes the
principal open problem of the paper (Section~\ref{P1:sec:open}).
In the thin-torus limit (restricting to the one-parameter ansatz family
$f=2\arctan(C/d)$), the companion paper~\cite{Paper2} proves the
fixed-point value exactly via Beta-function integration, yielding a
distinct fixed point $C^*\approx 1.817$ for that parametrisation; the two
values correspond to different approximations of the full toroidal geometry
and are not directly comparable.

%══════════════════════════════════════════════════════════════════════════════
\section{Numerical Results}
\label{P1:sec:numerics}
%══════════════════════════════════════════════════════════════════════════════

\subsection{Converged saddle-point solution}

All results were obtained by gradient descent on the feedback energy
$\Efb$ in the $Q=2$ topological sector, starting from the converged
standard FN Hopfion at $\bst=0$ and slowly increasing $\bst$ to $0.45225$
in steps of $0.01$--$0.05$, warm-starting each step.  The final run used a
$256\times256$ grid with $h=0.12$, $R_0=3.0$.  Code and the converged
profile array are deposited at~\cite{Code2026}.
The saddle-point profile is shown in Figures~\ref{P1:fig:profile}
and~\ref{P1:fig:profile_sym}; the two figures use different axis conventions
(half-plane vs.\ symmetric) and complement each other as described in their
captions.

\begin{center}
\begin{tabular}{llll}
\toprule
Quantity & Numerical value & Target (exact) & Gap \\
\midrule
$\Jfb/\Ja$ & $1.618030$ & $\vp = 1.618034$ & $-0.00025\%$ \\
$J_4/J_{2a}$ & $0.227168$ & $2^{4/3}/\vp^5 = 0.227214$ & $-0.020\%$ \\
$\sopt$ & $1.122350$ & $2^{1/6} = 1.122462$ & $-0.010\%$ \\
$\bz\rCMB$ & $10.012$ & $10$ (exact) & $+0.12\%$ \\
Virial sum & $1.99992$ & $2$ (exact) & $-0.004\%$ \\
\bottomrule
\end{tabular}
\end{center}

All five gaps are consistent in magnitude with a single $O(h^2)$
discretisation error at $h=0.12$.

\paragraph{Two-stage convergence structure.}
The numerical procedure has a natural two-stage structure that
illuminates the roles of shape and virial balance separately.
In the first stage, gradient descent on the
\emph{scale-invariant geometric mean functional}
$E_{\mathrm{geom}} = (\Ja + \ms\,\Jiso)\cdot\Jfour$
(Section~\ref{P1:sec:egeom}), starting from the rational BS ansatz,
converges to the Hopfion saddle shape in approximately
$10^4$ gradient steps. This profile already satisfies 
$\Jfour/\Ja \approx 0.2175$
($95.7\%$ of the WZW target $2^{4/3}/\vp^5$) and provides
$s_{\mathrm{opt}}\approx 1$ to within $0.05\%$ without any
feedback.  In the second stage, the feedback solver warm-starts from this profile
and scans $\bst$ until $\Jfb/\Ja = \vp$ is achieved,
refining $\Jfour/\Ja$ to $0.2272$ ($0.020\%$ gap).
At $\bst = 0.452$, the feedback solver converged in
$\approx 10^4$ additional steps. The first stage also provides the natural initial condition
for universality tests at other values of $R_0$: running
$E_{\mathrm{geom}}$ for $10^4$ steps at the target $R_0$ yields a
correctly shaped profile in minutes, avoiding the instabilities
that arise from translating the $R_0=3$ converged profile.

\subsection{The $R_0$-universality conjecture}
\label{P1:sec:universality}

The WZW algebra prediction~\eqref{P1:eq:general_formula} makes no reference
to the torus major radius $R_0$: the ratio $\Jfour/\Ja = 2^{4/3}/\vp^5$
is determined entirely by the level $k=3$ and the fusion count $N_{\mathrm{fus}}=2$,
both topological invariants of the $\mathrm{SU}(2)_3$ WZW model.
This implies the following conjecture.

\begin{quote}
\textbf{$R_0$-Universality Conjecture.}
At the self-consistent EL minimum of $\Efb$ in the $Q=2$ sector
(i.e., at the unique $\bst(R_0)$ for which $\Jfb/\Ja=\vp$),
the ratio $\Jfour/\Ja = 2^{4/3}/\vp^5$ is independent of the torus
major radius $R_0$.
\end{quote}

\paragraph{Numerical evidence.}
The conjecture is confirmed at $R_0=3$ to $0.020\%$ by the
converged feedback run of Section~\ref{P1:sec:numerics}.
Dedicated runs at $R_0=4$ and $R_0=5$ were performed using
the two-stage method of Section~\ref{P1:sec:numerics}: the
$E_{\mathrm{geom}}$ saddle-snapshot at each target $R_0$
was used as the initial condition for the feedback solver
(Section~\ref{P1:sec:egeom}).
At both $R_0=4$ and $R_0=5$ the feedback gradient flow found
no direction to move from the $E_{\mathrm{geom}}$ saddle:
the profile shape was unchanged across multiple solver sessions
(three sessions at $R_0=4$, two at $R_0=5$), confirming that
the $E_{\mathrm{geom}}$ saddle \emph{is} the $E_{\mathrm{fb}}$
saddle at larger $R_0$.
This is consistent with the self-similarity argument:
as $R_0$ increases the Hopf tube moves away from the $r=0$
axis, the cylindrical correction $1/r^2$ in $A = D_0^{-1}+r^{-2}$
becomes negligible at the tube, and $E_{\mathrm{geom}}$ captures
the tube geometry with increasing fidelity.

The resulting three-point dataset is:

\begin{center}
\renewcommand{\arraystretch}{1.2}
\begin{tabular}{ccccc}
\toprule
$R_0$ & $\Jfour/\Ja$ & ratio to $2^{4/3}/\vp^5$ & $\bst$ & Method \\
\midrule
$3$ & $0.22717$ & $0.9998$ & $0.452$ & Feedback converged ($24$ sessions) \\
$4$ & $0.21744$ & $0.9570$ & $0.483$ & $E_{\mathrm{geom}}$ saddle $=$ $E_{\mathrm{fb}}$ saddle \\
$5$ & $0.23160$ & $1.0193$ & $1.154$ & $E_{\mathrm{geom}}$ saddle $=$ $E_{\mathrm{fb}}$ saddle \\
\midrule
$\infty$ & $0.2301$ & $1.013$ & --- & $1/R_0$ extrapolation \\
\bottomrule
\end{tabular}
\end{center}

The $R_0=4$ and $R_0=5$ values bracket the WZW target from below
and above ($-4.3\%$ and $+1.9\%$ respectively), with the target
near the centre.
A least-squares fit $\Jfour/\Ja = g_\infty + a/R_0$ to the three
data points gives $g_\infty = 0.2301$ and $a = -0.0179$,
yielding a continuum estimate $1.3\%$ above the WZW prediction.
The residuals are oscillatory rather than monotone, consistent
with grid discretisation effects rather than a systematic
physical deviation.
Note also that $\bst$ is \emph{not} universal: it increases
from $0.452$ to $0.483$ to $1.154$ as $R_0$ grows, reflecting
the fact that the feedback coupling must compensate for the
changing magnitude of $\mathrm{kern}$ near $r=0$ as the Hopf
tube moves outward.
The ratio $\Jfour/\Ja$, which is dominated by the tube region
where both $J_4$ and $J_{2a}$ are concentrated, is the
$R_0$-independent quantity predicted by the WZW algebra.

The $R_0=3$ result remains the primary numerical confirmation
of the conjecture ($0.020\%$ gap, $O(h^2)$ discretisation
error).  The three-point dataset provides additional evidence
consistent with universality in the continuum limit, though a
precision confirmation at $R_0\neq 3$ comparable to the $R_0=3$
result would require the full multi-session feedback convergence
at each $R_0$ (comparable in cost to the original computation).

\subsection{Why Richardson extrapolation is non-trivial}

The functional $\Efb = \Kfb\cdot\Jfour$ has the trivial vacuum $f=0$ as
its global minimum.  Unconstrained gradient descent collapses to $f=0$
unless started from a profile already close to the saddle.  Projecting a
converged solution from grid spacing $h_1$ to a finer grid $h_2$ changes
the ratio $\Jfb/\Ja$ at step~1 (from $\vp$ to approximately $1.24$ at
$h_2/h_1=10/12$), which is sufficient to put the trajectory in the basin of
attraction of $f=0$.  A full Richardson extrapolation therefore requires
either a constrained $Q=2$ optimiser or a much more conservative
$h$-refinement strategy.  The $0.020\%$ gap at $h=0.12$ is our best
numerical estimate; the analytic route (Theorem~\ref{P1:thm:modBPS} plus
the open conjecture) is the cleaner path to confirming $\Jfour/\Ja =
2^{4/3}/\vp^5$ exactly.

%══════════════════════════════════════════════════════════════════════════════
\section{Status of Proofs}
\label{P1:sec:open}
%══════════════════════════════════════════════════════════════════════════════

\begin{center}
\renewcommand{\arraystretch}{1.25}
\begin{tabular}{p{6.5cm}p{4.5cm}p{3cm}}
\toprule
Statement & Status & Reference \\
\midrule
$1+\ms\vp = 2\vp$ & \textcolor{darkgreen}{Exact} & Lemma~\ref{P1:lem:golden} \\
$\vp^{-1}+\ms = 2$ & \textcolor{darkgreen}{Exact} & Lemma~\ref{P1:lem:virial} \\
$\lam=\vp^6$ & \textcolor{darkgreen}{Derived: unique soln.\ $Q_\mathrm{num}=Q_\mathrm{gp}$} & Cor.~\ref{P1:cor:lambda} \\
$\Jfb$ from parent Lagrangian~\eqref{P1:eq:Lkin} restricted to $Q=2$ & \textcolor{darkgreen}{Exact} & Sec.~\ref{P1:sec:vacuum_thm} \\
$Q=2$ icosahedral Hopfion: unique min.\ action at $\Lcond$ & \textcolor{darkgreen}{Proved (parts i--iv exact; part v uses Cor.~\ref{P1:cor:lambda})} & Thm.~\ref{P1:thm:vacuum} \\
$\Kfb = 2\vp\Ja$ at $\Jfb/\Ja=\vp$ & \textcolor{darkgreen}{Exact} & Prop.~\ref{P1:prop:Kfb} \\
$P_{\Kfb}=2\vp P_{\Ja}$ & \textcolor{darkgreen}{Exact} & Prop.~\ref{P1:prop:PKfb} \\
$\sopt=2^{1/6}\Leftrightarrow\Jfour/\Ja=2^{4/3}/\vp^5$ & \textcolor{darkgreen}{Exact} & Thm.~\ref{P1:thm:scale} \\
$K=2/3$ from $\mathbb{Z}_3$ Fourier & \textcolor{darkgreen}{Exact} & Thm.~\ref{P1:thm:koide} \\
$\dim_q(\tfrac{1}{2};k{=}3) = \vp$ & \textcolor{darkgreen}{Exact} & Thm.~\ref{P1:thm:dimq} \\
$I_4/I_{2a} = C(2C{+}1)/((C^2{+}1)(3C{-}1))$ & \textcolor{darkgreen}{Exact} & Thm.~\ref{P1:thm:modBPS} \\
$I_4/I_{2a}|_{\mathrm{BPS}} = 3/4$ & \textcolor{darkgreen}{Exact} & Prop.~\ref{P1:prop:bps} \\
$\sopt^{2k}=N_{\mathrm{fus}}$ (Linking Scale Conj.) & \textcolor{darkgreen}{Proved in Paper~II~\cite{Paper2}} & Conj.~\ref{P1:conj:linking} \\
$\Jfour/\Ja = 2^{4/3}/\vp^5$ at EL minimum & \textcolor{darkgreen}{Proved (Paper~II~\cite{Paper2})} & Sec.~\ref{P1:sec:open} \\
$\sinW = 3/(8\vp)$ & \textcolor{darkgreen}{Proved (Cor.~\ref{P1:cor:weinberg_unconditional})} & Thm.~\ref{P1:thm:weinberg}, \ref{P1:thm:amplitude_ratio}, Cor.~\ref{P1:cor:weinberg_unconditional} \\
$y_e = e^{-25\pi/6}$ & \textcolor{darkgreen}{Proved (Paper~II~\cite{Paper2})} & Result~\ref{P1:res:yukawa} \\
$\theta_K = 132.550^\circ$ & \textcolor{darkgreen}{Proved (Paper~II~\cite{Paper2})} & Result~\ref{P1:res:koide_angle} \\
\bottomrule
\end{tabular}
\end{center}

\newpage
\paragraph{Resolution of the analytic gap --- Paper~II.}
All predictions of this paper flowed from a single field-theoretic
statement --- that the EL minimum of $\Efb$ in the $Q=2$ sector
at the self-consistent $\bst$ satisfies $\Jfour/\Ja = 2^{4/3}/\vp^5$
exactly (in the continuum limit $h\to0$) --- which was confirmed
numerically to $0.020\%$ but not proved analytically in this paper.
By Theorem~\ref{P1:thm:scale} and Proposition~\ref{P1:prop:Kfb}, this is
equivalent to the EL equation taking the BPS form
\begin{equation}
  \frac{\delta\Kfb}{\delta f} \;=\; -\leff\,\frac{\delta\Jfour}{\delta f},
  \qquad
  \leff \;:=\; \frac{\Kfb}{\Jfour} \;=\; \frac{\lam}{2^{1/k}}\bigg|_{k=3}
  \;=\; \frac{\lam}{2^{1/3}},
  \label{P1:eq:BPS_EL}
\end{equation}
which is a special case of the \emph{level-$k$ general formula}
\begin{equation}
  \frac{\Jfour}{\Ja}\bigg|_{\mathrm{EL}}
  \;=\; \frac{2^{1+1/k}}{[\dim_q(\tfrac{1}{2};k)]^{2k-1}},
  \label{P1:eq:general_formula}
\end{equation}
confirmed numerically to $0.020\%$.
At $k=2$,~\eqref{P1:eq:general_formula} recovers $\Jfour/\Ja=1$,
the standard FN BPS condition in BS units, providing a non-trivial
consistency check.
The $k=3$ case is algebraically distinguished: the exponent $6-2k=0$
causes the $\dim_q$ factor to cancel, leaving $\sopt = 2^{1/6}$
independent of the quantum-group deformation parameter $\vp$.

Numerically, on the converged $256\times256$ profile at $h=0.12$:
\begin{equation}
  \frac{\Kfb}{\Jfour}\bigg|_{\mathrm{grid}}
  \;=\; 14.246
  \;\approx\; \frac{\lam}{2^{1/3}} = 14.242
  \quad (+0.03\%),
  \label{P1:eq:leff_numerical}
\end{equation}
consistent with $O(h^2)$ discretisation error at $h=0.12$.

The proof of~\eqref{P1:eq:general_formula} and~\eqref{P1:eq:BPS_EL}
requires three steps:
\begin{enumerate}[noitemsep,label=\textup{(\roman*)}]
  \item \textbf{Existence and uniqueness.}
    For each $\leff>0$, the BPS equation
    $\delta\Kfb/\delta f = -\leff\,\delta\Jfour/\delta f$
    has a unique smooth solution $f_{\leff}$ in the $Q=2$ sector.
  \item \textbf{Monotonicity.}
    The map $H\colon \leff \mapsto \Kfb[f_{\leff}]/\Jfour[f_{\leff}]$
    is strictly monotone, so the fixed-point equation $H(\leff^*)=\leff^*$
    has a unique solution.
  \item \textbf{Fixed-point value.}
    The unique fixed point satisfies $\leff^* = \lam/2^{1/3}$,
    equivalently $\Jfour/\Ja = 2^{4/3}/\vp^5$.
    Given the proved identity $\Kfb = 2\vp\Ja$ (Proposition~\ref{P1:prop:Kfb}),
    this reduces to a single algebraic statement:
    \begin{equation}
      \sopt^{2k}\big|_{k=3} \;=\; N_{\mathrm{fus}} \;=\; 2,
      \label{P1:eq:linking_scale}
    \end{equation}
    where $N_{\mathrm{fus}} = |\{\text{fusion channels of }\phi_{1/2}
    \text{ at level }k\}|= 2$ is proved (the rule
    $\phi_{1/2}\times\phi_{1/2}=\phi_0+\phi_1$ at $k=3$).
    Equation~\eqref{P1:eq:linking_scale} is the
    \emph{Linking Scale Conjecture}; its precise form
    is stated below (Conjecture~\ref{P1:conj:linking}).
\end{enumerate}
Steps~(i) and~(ii) are standard arguments established in
Paper~II~\cite{Paper2}.
Step~(iii) is proved in Paper~II~\cite{Paper2} via a WZW Ward
identity showing that $\Kfb/\Jfour$ is constant on the EL manifold
with value $\lam/N_{\mathrm{fus}}^{1/k}$.
As a consequence, all results previously labelled ``conditional on
the $\Jfb/\Ja=\vp$ axiom'' --- including $y_e=e^{-25\pi/6}$,
$\theta_K=132.550^\circ$, and the lepton mass ratios --- are
fully proved without any remaining assumption.

\begin{conjecture}[Linking Scale Conjecture --- proved in Paper~II~\cite{Paper2}]
\label{P1:conj:linking}
For the $Q$-Hopfion at WZW level $k$, the EL saddle of $\Efb$
in the topological sector with virial condition
$\Jfb/\Ja = \dim_q(\tfrac{1}{2};k)$ satisfies
\begin{equation}
  \sopt^{2k} \;=\; N_{\mathrm{fus}}(Q,k),
  \label{P1:eq:lsc}
\end{equation}
where $N_{\mathrm{fus}}(Q,k)$ is the number of fusion channels of
$\phi_{1/2}$ at level $k$.
Equivalently, the effective BPS coupling satisfies
$\leff^* = \lam/N_{\mathrm{fus}}^{1/k}$.
\end{conjecture}

\paragraph{Verification at $k=2$ and $k=3$.}
At $k=2$ (standard FN): $N_{\mathrm{fus}}=1$
($\phi_{1/2}\times\phi_{1/2}=\phi_0$ only at level $2$),
so~\eqref{P1:eq:lsc} gives $\sopt^4=1$, equivalently
$\Jfour/\Ja=1$, which is the standard FN BPS condition
in BS units.  This is exact.
At $k=3$: $N_{\mathrm{fus}}=2$, so~\eqref{P1:eq:lsc} gives
$\sopt^6=2$, which is equivalent to
$\Jfour/\Ja = 2^{4/3}/\vp^5$ by pure algebra
(using $\Kfb=2\vp\Ja$ and $\lam=\vp^6$).
This is confirmed numerically to $0.020\%$ at $R_0=3$
and is consistent with the three-point
$R_0\in\{3,4,5\}$ dataset of Section~\ref{P1:sec:universality}.

\paragraph{Proof strategy.}
The conjecture cannot follow from the EL equation alone:
the BPS form $\delta\Kfb/\delta f = -\leff\,\delta\Jfour/\delta f$
holds for any $\leff$ and does not fix its value.
The value is pinned by the topology of the $Q=2$ sector.
The proof requires establishing a \emph{WZW Ward identity}
for $\Efb$: that $\Kfb/\Jfour$ is constant on the EL
manifold and equal to $\lam/N_{\mathrm{fus}}^{1/k}$ by the
WZW character of the topological sector.
The virial condition $\Jfb/\Ja = \vp = \dim_q(\tfrac{1}{2};k{=}3)$
is the field-theoretic avatar of the WZW $S$-matrix
self-duality $S_{1/2,\,1/2}=S_{0,0}$ proved in
Theorem~\ref{P1:thm:dimq}; establishing the corresponding
Ward identity for $\Kfb/\Jfour$ would complete the proof.
The proof is given in Paper~II~\cite{Paper2}.

%══════════════════════════════════════════════════════════════════════════════
\section{Summary and Complete Prediction Table}
\label{P1:sec:summary}
%══════════════════════════════════════════════════════════════════════════════

\begin{center}
\renewcommand{\arraystretch}{1.3}
\begin{tabular}{p{4cm}p{5cm}p{1.5cm}p{1.5cm}c}
\toprule
Quantity & Formula & Predicted & Observed & Deviation \\
\midrule
\multicolumn{5}{l}{\textit{Exact algebraic identities}} \\
$1+\ms\vp$ & $= 2\vp$ & exact & --- & --- \\
$\vp^{-1}+\ms$ & $= 2$ & exact & --- & --- \\
$\ms$ & $= 3-\vp = 1+\vp^{-2}$ & $1.38197$ & confirmed & $0.08\%$ \\
$\lam=\vp^6$ & unique soln.\ $Q_\mathrm{num}=Q_\mathrm{gp}$ & $17.944$ & confirmed & --- \\
$Q=\ord(q_{2I})$ & $= 10$ & exact & $10.012$ (lattice) & $0.12\%$ \\
$K$ (Koide) & $= 2/3$ from $\mathbb{Z}_3$ & exact & $0.66666$ & $<0.01\%$ \\
$I_4/I_{2a}|_{\mathrm{BPS}}$ & $= 3/4$ & exact & $0.750\pm0.001$ & $<0.01\%$ \\[4pt]
\multicolumn{5}{l}{\textit{Proved predictions (Linking Scale Conjecture resolved in Paper~II~\cite{Paper2})}} \\
$\sinW$ & $3/(8\vp)$ & $0.23176$ & $0.23122$ & $0.23\%$ \\
$M_W/M_Z$ & $\sqrt{1-3/(8\vp)}$ & $0.8765$ & $0.8815$ & $0.56\%$ \\
$y_e$ & $e^{-25\pi/6}$ & $2.066\times10^{-6}$ & $2.077\times10^{-6}$ & $0.46\%$ \\
$\vEW$ & $m_e e^{25\pi/6}$ & $247.4\;\text{GeV}$ & $246.2\;\text{GeV}$ & $0.46\%$ \\
$\theta_K$ & WZW holonomy & $132.550^\circ$ & $132.732^\circ$ & $0.14\%$ \\
$m_\mu/m_e$ & Koide+WZW & $206.69$ & $206.77$ & $0.04\%$ \\
$m_\tau/m_e$ & Koide+WZW & $3476.2$ & $3477.2$ & $0.03\%$ \\
\bottomrule
\end{tabular}
\smallskip
\end{center}

\bigskip

\paragraph{Inputs used.}
The only numerical input taken from experiment is $\TCMB=2.7255\;\text{K}$
(fixes $\bz$ and $\rCMB$).  The electron mass $m_e$ is \emph{derived}:
Paper~XII~\cite{Paper12} proves the one-parameter chain
$\TCMB\to\rCMB=10\to\vp^9\sqrt{\Ja\Jfour}=10\to m_e$ (residual $0.013\%$
after the $\overline{\mathrm{MS}}$-to-pole scheme correction $-\alpha/\pi$
of Paper~IV~\cite{Paper4}).
All other SM quantities are predicted.  The golden ratio $\vp$ 
enters through two independent proved routes:
as the quantum dimension $\dim_q(\tfrac{1}{2};k{=}3)$ of the
$\mathrm{SU}(2)_3$ WZW model (Theorem~\ref{P1:thm:dimq}), and as the
character $\chi_{\Gamma_2}(C_5)$ of the fundamental $2I$ representation
(Section~\ref{P1:sec:group}).
The Bogomolny parameter $\lam=\vp^6$ is the unique value consistent with
$Q_{\mathrm{num}}=Q_{\mathrm{group}}=10$ (Corollary~\ref{P1:cor:lambda}).
The Weinberg angle $\sinW=3/(8\vp)$ factorises into two irreducible
components, $3/8$ from the $\mathrm{SU}(5)$ kinematic trace and $1/\vp$
from the WZW $S$-matrix; both components are fully proved ---
the amplitude ratio $\mathcal{A}_Y/\mathcal{A}_{T_3}=1/\vp$ follows
from the geometry of the Hopf submersion
(Theorem~\ref{P1:thm:amplitude_ratio}), making the Weinberg angle
unconditional (Corollary~\ref{P1:cor:weinberg_unconditional}).
The self-consistency condition $\Jfb/\Ja=\vp$ is confirmed numerically
to $0.00025\%$; its analytic derivation reduces to the Linking Scale
Conjecture (Conjecture~\ref{P1:conj:linking}), proved in
Paper~II~\cite{Paper2} (Section~\ref{P1:sec:open}). The quark Koide relation $K_{\mathrm{up}}=K_{\mathrm{down}}=2/3$ is proved
exactly at the condensate scale $\LUV$ in Paper~V~\cite{Paper5}
(Theorem~7.4 therein), as an exact consequence of the universal
colour level shift $\delta n^{(C)}=1/2$ and the lepton Koide
theorem above.  The level assignments for the quark $\mathrm{SU}(3)_C$
representations are fixed by the $\mathbb{CP}^2$ Hopf fibration
(Paper~V, Section~3), confirming the prediction of this section.
The three CKM mixing angles and the CP-violating phase $\delta_{\rm CP}$
are likewise derived in Paper~V with no free parameters
(Theorems~8.4 and~8.7 therein).

\paragraph{Open problems ranked by priority.}
\begin{enumerate}[noitemsep]
  \item \textbf{[Resolved in Paper~II~\cite{Paper2}]} Prove the Linking Scale Conjecture
    (Conjecture~\ref{P1:conj:linking}): that the EL saddle of $\Efb$
    in the $Q=2$ topological sector satisfies $\sopt^{2k}=N_{\mathrm{fus}}$,
    equivalently $\Jfour/\Ja = 2^{4/3}/\vp^5$.
    Given the proved identity $\Kfb=2\vp\Ja$, this reduces to
    establishing a single WZW Ward identity: that $\Kfb/\Jfour$
    is constant on the EL manifold and equal to $\lam/N_{\mathrm{fus}}^{1/k}$.
    The conjecture is verified at $k=2$ (standard FN, exact) and
    $k=3$ (numerically to $0.020\%$; consistent with the
    three-point $R_0\in\{3,4,5\}$ dataset).
    The proof is given in Paper~II~\cite{Paper2}.
    \item \textbf{[Resolved in Paper~V~\cite{Paper5}]}
      Quark-sector Koide from the Hopf tower.
      The relation $K_{\mathrm{up}}=K_{\mathrm{down}}=2/3$ is proved exactly
      at the condensate scale (Theorem~7.4 of Paper~V) as a consequence of
      the universal colour level shift $\delta n^{(C)}=1/2$ from the
      $\mathbb{CP}^2$ Hopf fibration: since every quark mass is a
      generation-universal rescaling of the corresponding lepton mass
      ($m_{u,c,t}^{(g)}=\vp^{-1}m_{\ell}^{(g)}$,
       $m_{d,s,b}^{(g)}=\vp^{-3}m_{\ell}^{(g)}$),
      and the Koide ratio is scale-invariant
      (Lemma~7.3 of Paper~V), the quark Koide ratio inherits the value
      $K=2/3$ from the \emph{lepton} Koide theorem proved in this paper
      (Theorem~7.4 herein).
      The two Koide results therefore have different origins:
      the lepton $K=2/3$ follows from the $\mathbb{Z}_3$ symmetry of the
      WZW T-matrix phases $\{T_1,T_2,T_3\}$ under $g\to g+1\pmod{3}$
      (a property of this paper's WZW conformal data);
      the quark $K=2/3$ is an \emph{inheritance}, exact because the
      colour shift is universal across all three generations.
      Paper~VI~\cite{Paper6} (Proposition~5.5 and Remark~5.6 therein)
      further proves that individual lepton tower levels are irrational
      (from the transcendental incommensurability of $\pi$ and $\ln\vp$),
      while quark-to-lepton \emph{ratios} are integer-exponent exact
      precisely because the $\ln\vp$ terms cancel between numerator and
      denominator --- the structural reason the two sectors behave
      differently in the tower.
  \item \textbf{[Resolved in Paper~V~\cite{Paper5}]}
    CKM/PMNS mixing angles from Coxeter exponents of $E_6,E_7,E_8$.
    The three CKM angles $|V_{us}|=\sqrt{m_d/m_s}$,
    $|V_{cb}|=e^{-\pi}$, $|V_{ub}|=\sqrt{m_u/m_t}$ and
    the CP-violating phase $\delta_{\rm CP}=65.0^\circ$
    are derived with no free parameters via the $2I$-McKay--WZW spectrum.
  \item \textbf{[Resolved in Papers~IV and~VI]} Derive $\vEW$ from $\TCMB$
    and $\vp$ alone (close the remaining two-scale gap).
    \emph{Resolved.}
    Paper~IV~\cite{Paper4} derives the exact formula
    \begin{equation*}
      \vEW \;=\; \Lcond\cdot\vp^{20}\cdot
      \exp\!\Bigl(\tfrac{49\pi}{6} - \tfrac{3}{400}\Bigr)
      \;\approx\; 247.9\,\mathrm{GeV},
    \end{equation*}
    accurate to $0.69\%$ (residual is $O(\alpha_s^2)$), with no Standard
    Model parameters.
    The three factors are: $\vp^{20}$ from the bosonic tower at level
    $Q=10$ (topological charge, proved in this paper); $e^{49\pi/6}$
    from the Chern--Simons action ($e^{4\pi}$, Paper~IV~\cite{Paper4}) combined with
    the electron Yukawa $y_e=e^{-25\pi/6}$ (proved in this paper); and
    $e^{-3/400}$ from the modular $T$-matrix correction
    $T_{1/2}/Q = 3/400$ (Paper~IV, exact WZW datum).
    The tower interpretation --- $\vp^{20}$ as the $n=-Q$ level of the
    bosonic tower, $e^{49\pi/6}$ as the Chern--Simons spoke amplitude ---
    is given in Paper~VI~\cite{Paper6} (Theorem~6.1 therein).
    No free parameters remain.
    Paper~XII~\cite{Paper12} proves $\vp^9\sqrt{\Ja\Jfour}=10$
    from $\TCMB$ alone via the CMB thermodynamic equilibrium condition,
    completing the single-input derivation:
    $\TCMB\to\Lcond\to\rho_{\rm CMB}=10\to\vp^9\sqrt{\Ja\Jfour}=10
    \to m_e=\Lcond\,\vp^{20}\,e^{4\pi-3/400-\alpha/\pi}$
    (residual $0.013\%$; the three exponent terms are: $4\pi$ from
    the Chern--Simons spoke, $-3/400$ from the WZW $T$-matrix phase,
    and $-\alpha/\pi$ from the one-loop $\overline{\mathrm{MS}}$-to-pole
    scheme conversion of Paper~IV~\cite{Paper4}).
\end{enumerate}

\newpage

%══════════════════════════════════════════════════════════════════════════════
\begin{thebibliography}{99}
%══════════════════════════════════════════════════════════════════════════════

\bibitem{Paper2}
F.~Manfredi,
\textit{BPS Structure and Fixed-Point Theorem for the
  Density-Feedback Faddeev--Niemi Hopfion},
Zenodo (2026).
\href{https://doi.org/10.5281/zenodo.19363491}{doi:10.5281/zenodo.19363491}

\bibitem{Paper3}
F.~Manfredi,
\textit{Two Constants from One Knot: The Fine Structure Constant and
  the Linking Scale as Exact Predictions of the Icosahedral WZW Condensate},
Zenodo (2026).
\href{https://doi.org/10.5281/zenodo.19478629}{doi:10.5281/zenodo.19478629}

\bibitem{Paper4}
F.~Manfredi,
\textit{The Hopf Spoke, WZW Fermion Mass Renormalization,
  and Three- and Four-Term Formulas for the Fine Structure Constant},
Preprint (2026).
\href{https://doi.org/10.5281/zenodo.19504729}{doi:10.5281/zenodo.19504729}

\bibitem{Paper5}
F.~Manfredi,
\textit{Colour Confinement, the QCD Scale, and Hadronic Structure
  from the Density-Feedback Hopfion},
Preprint (2026).
\href{https://doi.org/10.5281/zenodo.20210460}{doi:10.5281/zenodo.20210460}

\bibitem{Paper6}
F.~Manfredi,
\textit{The Golden-Ratio Tower: Fermion Scale Hierarchy from the Hopfion RG Fixed Point},
Preprint (2026).
\href{https://doi.org/10.5281/zenodo.19646945}{doi:10.5281/zenodo.19646945}

\bibitem{Paper12}
F.~Manfredi,
\textit{Paper~XII: The Profile Normalisation Conjecture --- Proof via CMB Energy Identity and Two-Route Equivalence},
\href{https://doi.org/10.5281/zenodo.19646945}{doi:10.5281/zenodo.19646945}

\bibitem{Code2026}
F. Manfredi,
\textit{Code and converged profile for the density-feedback
  Faddeev--Niemi Hopfion} ($256\times256$, $\beta^*=0.452$),
Zenodo, (2026) \href{https://doi.org/10.5281/zenodo.19340941}{doi:10.5281/zenodo.19340941}

\bibitem{Faddeev1997}
L.~D. Faddeev and A.~J. Niemi,
\textit{Stable knot-like structures in classical field theory},
Nature \textbf{387}, 58 (1997).
\href{https://doi.org/10.1038/387058a0}{doi:10.1038/387058a0}

\bibitem{Battye1998}
R.~A. Battye and P.~M. Sutcliffe,
\textit{Knots as stable soliton solutions in a three-dimensional classical
field theory},
Phys.\ Rev.\ Lett.\ \textbf{81}, 4798 (1998).
\href{https://doi.org/10.1103/PhysRevLett.81.4798}{doi:10.1103/PhysRevLett.81.4798}

\bibitem{McKay1980}
J.~McKay,
\textit{Graphs, singularities, and finite groups},
Proc.\ Symp.\ Pure Math.\ \textbf{37}, 183 (1980).
\href{https://doi.org/10.1090/pspum/037}{doi:10.1090/pspum/037}

\bibitem{Verlinde1988}
E.~Verlinde,
\textit{Fusion rules and modular transformations in 2D conformal field theory},
Nucl.\ Phys.\ B \textbf{300}, 360 (1988).
\href{https://doi.org/10.1016/0550-3213(88)90603-7}{doi:10.1016/0550-3213(88)90603-7}

\bibitem{Adams1960}
J.~F. Adams,
\textit{On the non-existence of elements of Hopf invariant one},
Ann.\ Math.\ \textbf{72}, 20 (1960).
\href{https://doi.org/10.2307/1970147}{doi:10.2307/1970147}

\bibitem{Hurwitz1898}
A.~Hurwitz,
\textit{\"Uber die Composition der quadratischen Formen von beliebig vielen
Variablen},
Nachr.\ Ges.\ Wiss.\ G\"ottingen (1898), 309.
ISBN 978-3-0348-4085-9.
% No DOI available for this 19th-century proceedings article.

\bibitem{Koide1983}
Y.~Koide,
\textit{A fermion-boson composite model of quarks and leptons},
Phys.\ Lett.\ B \textbf{120}, 161 (1983).
\href{https://doi.org/10.1016/0370-2693(83)90644-5}{doi:10.1016/0370-2693(83)90644-5}

\bibitem{KhouryWeltman2004}
J.~Khoury and A.~Weltman,
\textit{Chameleon fields: Awaiting surprises for tests of gravity in space},
Phys.\ Rev.\ Lett.\ \textbf{93}, 171104 (2004).
\href{https://doi.org/10.1103/PhysRevLett.93.171104}{doi:10.1103/PhysRevLett.93.171104}

\bibitem{White2026}
H.~White, J.~Vera, A.~Sylvester, and L.~Dudzinski,
\textit{Emergent quantization from a dynamic vacuum},
Phys.\ Rev.\ Research \textbf{8}, 013264 (2026).
\href{https://doi.org/10.1103/l8y7-r3rm}{doi:10.1103/l8y7-r3rm}

\bibitem{PDG2024}
R.~L. Workman et al.\ (Particle Data Group),
\textit{Review of Particle Physics},
Prog.\ Theor.\ Exp.\ Phys.\ \textbf{2022}, 083C01 (2022);
updated online 2024.
\href{https://doi.org/10.1093/ptep/ptac097}{doi:10.1093/ptep/ptac097}
%Journal

\bibitem{Derrick1964}
G.~H. Derrick,
\textit{Comments on nonlinear wave equations as models for elementary particles},
J.\ Math.\ Phys.\ \textbf{5}, 1252 (1964).
\href{https://doi.org/10.1063/1.1704233}{doi:10.1063/1.1704233}

\bibitem{BornInfeld1934}
M.~Born and L.~Infeld,
\textit{Foundations of the new field theory},
Proc.\ R.\ Soc.\ London A \textbf{144}, 425 (1934).
\href{https://doi.org/10.1098/rspa.1934.0059}{doi:10.1098/rspa.1934.0059}

\bibitem{Madelung1927}
E.~Madelung,
\textit{Quantentheorie in hydrodynamischer Form},
Z.\ Phys.\ \textbf{40}, 322 (1927).
\href{https://doi.org/10.1007/BF01400372}{doi:10.1007/BF01400372}

\bibitem{VK1979}
A.~F. Vakulenko and L.~V. Kapitanski,
\textit{Stability of solitons in $S^2$ in the nonlinear $\sigma$-model},
Dokl.\ Akad.\ Nauk SSSR \textbf{246}, 840--842 (1979).
% No DOI available for this Soviet-era Doklady article.

\bibitem{GeorgiGlashow1974}
H.~Georgi and S.~L. Glashow,
\textit{Unity of all elementary-particle forces},
Phys.\ Rev.\ Lett.\ \textbf{32}, 438 (1974).
\href{https://doi.org/10.1103/PhysRevLett.32.438}{doi:10.1103/PhysRevLett.32.438}

\bibitem{Cardy1989}
J.~L. Cardy,
\textit{Boundary conditions, fusion rules and the Verlinde formula},
Nucl.\ Phys.\ B \textbf{324}, 581 (1989).
\href{https://doi.org/10.1016/0550-3213(89)90521-X}{doi:10.1016/0550-3213(89)90521-X}

\end{thebibliography}

\end{document}

\typeout{get arXiv to do 4 passes: Label(s) may have changed. Rerun}